% %%%%%%%%%%%%%%%%%%%%%%%%%%%%%%%%%%%%%%%%%%%%%%%%%%%%%%%%%%%%%%%%%%%%%%% %
%                                                                         %
% The Project Gutenberg EBook of How to Draw a Straight Line, by A.B. Kempe
%                                                                         %
% This eBook is for the use of anyone anywhere at no cost and with        %
% almost no restrictions whatsoever.  You may copy it, give it away or    %
% re-use it under the terms of the Project Gutenberg License included     %
% with this eBook or online at www.gutenberg.org                          %
%                                                                         %
%                                                                         %
% Title: How to Draw a Straight Line                                      %
%        A Lecture on Linkages                                            %
%                                                                         %
% Author: A.B. Kempe                                                      %
%                                                                         %
% Release Date: April 24, 2008 [EBook #25155]                             %
%                                                                         %
% Language: English                                                       %
%                                                                         %
% Character set encoding: ISO-8859-1                                      %
%                                                                         %
% *** START OF THIS PROJECT GUTENBERG EBOOK HOW TO DRAW A STRAIGHT LINE ***
%                                                                         %
% %%%%%%%%%%%%%%%%%%%%%%%%%%%%%%%%%%%%%%%%%%%%%%%%%%%%%%%%%%%%%%%%%%%%%%% %

\def\ebook{25155}
\newtoks\PGheader
{\catcode`\#11\relax\catcode`\L\active\obeylines\obeyspaces%
\global\PGheader={%
The Project Gutenberg EBook of How to Draw a Straight Line, by A.B. Kempe

This eBook is for the use of anyone anywhere at no cost and with
almost no restrictions whatsoever.  You may copy it, give it away or
re-use it under the terms of the Project Gutenberg License included
with this eBook or online at www.gutenberg.org


Title: How to Draw a Straight Line
       A Lecture on Linkages

Author: A.B. Kempe

Release Date: April 24, 2008 [EBook #25155]

Language: English

Character set encoding: ISO-8859-1

*** START OF THIS PROJECT GUTENBERG EBOOK HOW TO DRAW A STRAIGHT LINE ***
}}
\AtBeginDocument{\CreditsLine{%
Produced by Joshua Hutchinson, David Wilson and the Online
Distributed Proofreading Team at http://www.pgdp.net (This
file was produced from images from the Cornell University
Library: Historical Mathematics Monographs collection.)
}}

%%%%%%%%%%%%%%%%%%%%%%%%%%%%%%%%%%%%%%%%%%%%%%%%%%%%%%%%%%%%%%%%%%%%%%%%%%%
%%                                                                       %%
%% Packages and substitutions:                                           %%
%%                                                                       %%
%% memoir:   Advanced book class. Required.                              %%
%% memhfixc: Part of memoir; needed to work with hyperref. Required.     %%
%% amsmath:  AMS mathematics enhancements. Required.                     %%
%% graphicx: Standard interface for graphics inclusion. Required.        %%
%%           Driver option needs to be set explicitly.                   %%
%% hyperref: Hypertext embellishments for pdf output. Required.          %%
%%           Driver option needs to be set explicitly.                   %%
%% indentfirst: Standard package to indent first line following chapter/ %%
%%              section headings. Recommended.                           %%
%% soul:     Facilitates effects like letterspacing. Recommended.        %%
%% wrapfig:  Allows placement of graphics inside text cutouts.           %%
%%             Strongly recommended.                                     %%
%%                                                                       %%
%%                                                                       %%
%% Producer's Comments: A fairly straightforward text, except for        %%
%%                      keeping the illustrations in sequence and not    %%
%%                      too far from the text referring to them.         %%
%%                                                                       %%
%% Things to Check:                                                      %%
%%                                                                       %%
%% Figure 23 fits snugly at the bottom of page 30 (pdf page 38): OK      %%
%% hyperref and graphicx driver option matches workflow: OK              %%
%% color driver option matches workflow (color package is called         %%
%%    by hyperref, so may rely on color.cfg): OK                         %%
%% Spellcheck: OK                                                        %%
%% Smoothreading pool: No                                                %%
%% LaCheck: OK                                                           %%
%% Lprep/gutcheck: OK                                                    %%
%% PDF pages:  63                                                        %%
%% PDF page size: 422 x 652pt (non standard)                             %%
%% PDF bookmarks: created (point to figures) but closed by default       %%
%% PDF document info: filled in                                          %%
%% PDF Reader displays document title in window title bar                %%
%% Images: 35 png and 1 pdf                                              %%
%% No overfull boxes, one underfull hbox and ten underfull vboxes        %%
%%    (caused by placement of images)                                    %%
%%                                                                       %%
%% pdflatex x2                                                           %%
%%                                                                       %%
%% Compile History:                                                      %%
%%                                                                       %%
%% Apr 08: dcwilson.                                                     %%
%%         Compiled with pdfLaTeX TWO times.                             %%
%%         MiKTeX 2.7, Windows XP Pro                                    %%
%%                                                                       %%
%%                                                                       %%
%%%%%%%%%%%%%%%%%%%%%%%%%%%%%%%%%%%%%%%%%%%%%%%%%%%%%%%%%%%%%%%%%%%%%%%%%%%

\listfiles
\documentclass[12pt,twoside,openright,onecolumn]{memoir}[2005/09/25]

% the page size is chosen to roughly approximate the line length
% and the number of lines per page of the original scans
\setstocksize{230mm}{149mm}
\settrimmedsize{\stockheight}{\stockwidth}{*}
\setlrmarginsandblock{1.8cm}{2.1cm}{*}
\setulmarginsandblock{2.1cm}{2.5cm}{*}
\setlength{\headsep}{14pt}
\setlength{\footskip}{0.6cm}
\fixthelayout
\typeoutlayout

\AtBeginDocument{\baselineskip=1.15\baselineskip} % increase leading like the original book

\makeatletter

% font issues
% Courier, for the PG licence stuff
\DeclareRobustCommand\ttfamily % Courier, for the PG licence stuff
        {\not@math@alphabet\ttfamily\mathtt
         \fontfamily{pcr}\fontencoding{T1}\selectfont}
% CM Dunhill as the only "standard" condensed font available
\DeclareFontFamily{OT1}{cmdh}{\hyphenchar\font45 }
\DeclareFontShape{OT1}{cmdh}{m}{n}{%
      <10><10.95><12><14.4><17.28><20.74><24.88>cmdunh10%
    }{}
\newcommand*\cmdh{\fontfamily{cmdh}\fontencoding{OT1}\selectfont}

\usepackage[latin1]{inputenc}[2006/05/05]
\usepackage{amsmath}[2000/07/18]
\IfFileExists{soul.sty}{% load soul package for letterspaced headings
  \usepackage{soul}[2003/11/17]
  }{% else make \textso a no-op
  \providecommand\textso[1]{##1}
  }
\usepackage{indentfirst}[1995/11/23] % to override default behaviour

% half-title and copyright pages
\let\transcribersnotes\@empty
\let\transcribersNotes\@empty
\newcommand{\transcribersnote}[1]{%
  \@ifnotempty{#1}{\g@addto@macro\transcribersnotes{#1\par}%
    \@xp\@ifempty\@xp{\transcribersNotes}%
      {\renewcommand{\transcribersNotes}{Transcriber's note}}
      {\renewcommand{\transcribersNotes}{Transcriber's notes}}}}
\newcommand{\CreditsLine}[1]{\newcommand{\thePr@ductionTeam}{#1}}

\def\makehalftitlepage{% the boilerplate header followed by a real half-title page
  \begingroup
  \begin{PGboilerplate}
    \PGlicencelink
    \the\PGheader
  \end{PGboilerplate}
  \null\vfil
  \cleartorecto
  \null\vskip0ptplus1.5fil
  \includegraphics[scale=1.8]{images/illus0003.png}
  \bigskip
  \begin{center}
  \spaceskip=0.75em
  \large
  HOW TO DRAW A STRAIGHT LINE;\\[16pt]
  \tiny
  A\\[16pt]
  \normalsize
  LECTURE ON LINKAGES
  \end{center}
  \cleartorecto
  \endgroup}

\def\makecopyrightpage{% production credits and transcriber's notes
  \begingroup
  \null\vfil
  \begin{center}\small
    \thePr@ductionTeam
  \end{center}
  \vfil\vfil
  \vbox{\Small\hsize=.75\textwidth\parindent=\z@\parskip=.5em
  \textit{\transcribersNotes}\par\medskip\raggedright
  \transcribersnotes\par}
  \cleartorecto
  \endgroup}

% title page
\aliaspagestyle{title}{empty}
% couldn't find a font to match the scan, hence the mathcal T in the middle
\renewcommand\maketitlehooka{\vspace*{-4em}\noindent
  \setbox0=\hbox{\textit{\Large NA$\mathcal T$\kern-.15emURE SERIES.}}%
  \copy0\vspace*{.25em}\hrule width\wd0 depth0pt height.4pt\vspace*{4em}}
\pretitle{\begin{center}\cmdh\Large}
\posttitle{\par\end{center}}
\preauthor{\vfill\begin{center}{\tiny BY}\\[12pt]\large}
\postauthor{\par\end{center}}
\def\affiliation#1{\renewcommand{\maketitlehookc}{\begin{center}
  \@setfontsize\miniscule\@vipt{9}
  #1\par\end{center}}}
\def\subtitle{\def\SubTitle}
\predate{\vfill\begin{center}\scriptsize\so\SubTitle\par
 \vspace{\z@\@plus2fill}\small}
\postdate{\par\end{center}\vspace*{-2em}}

% headers and footers
\let\Small\footnotesize
\makepagestyle{kempe}
  \makeevenhead{kempe}{\Small\thepage}{\Small\slshape HOW TO DRAW A STRAIGHT LINE:}{}
  \makeoddhead{kempe}{}{\Small\slshape A LECTURE ON LINKAGES}{\Small\thepage}
  % make it hard to end the next (odd) page with a hyphen
  \makeevenfoot{kempe}{}{}{\global\brokenpenalty10000}
  % make it less hard to end the next (even) page with a hyphen
  % (because the other end of the hyphen will still be on the same spread)
  \makeoddfoot{kempe}{}{}{\global\brokenpenalty150}
\pagestyle{kempe}
\copypagestyle{chapter}{empty}
\copypagestyle{licence}{headings}
\makeevenhead
  {licence}{\normalfont\Small\thepage}{\normalfont\Small\slshape LICENSING}{}
\makeoddhead
  {licence}{}{\normalfont\Small\slshape LICENSING}{\normalfont\Small\thepage}

\brokenpenalty10000

\renewcommand{\chaptitlefont}{\cmdh\Large\centering}
\setlength{\beforechapskip}{40mm}
\setlength{\afterchapskip}{24pt}
\setsecheadstyle{\normalfont\cmdh\centering}
\setsubsecheadstyle{\normalfont\large\centering}
\def\Sectionformat#1#2{\textso{#1}} % subvert hyperref hook to letterspace section headings

% illustrations
% (external files are .png, plus one .pdf)
% defaults are not stretchy enough
\setlength\textfloatsep{20\p@ \@plus 6\p@ \@minus 4\p@}
\setlength\intextsep   {14\p@ \@plus 8\p@ \@minus 4\p@}
\renewcommand\bottomfraction{0.7}
% driver should be specified in graphics.cfg;
% if not, add explicit option to graphicx call
\usepackage[final]{graphicx}[1999/02/16]
\GenericWarning{*** }{***\MessageBreak
 Important Note: this document comes with PNG and PDF graphics,\MessageBreak
 so make sure you use an appropriate workflow!\MessageBreak
 \expandafter\@gobble\@gobble}
% to simplify making hyperlinks to illustrations
\newcommand\figref[2][Fig.]{\hyperlink{fig:#2}{#1~#2}}
% to insert an illustration, with an anchor at the top, and a bookmark entry
\newcommand\Figure[2]{\setbox0=\hbox to\linewidth
  {\hss\includegraphics[scale=1.77]{images/illus0#2.png}%
  \hss}\leavevmode\llap{\Hy@writebookmark{}{Figure #1}{fig:#1}{1}{toc}}%
  \raise\ht0\hbox{\DPanchor{fig:#1}}\box0}
% to wrap text around Fig. 23
\IfFileExists{wrapfig.sty}{% load package
  \RequirePackage{wrapfig}[2003/01/31]
  }{ % else make wrapfigure effectively just a normal figure*
  \def\wrapfigure[##1]##2##3##4\end##5{%
  \begin{figure*}[hbt]##4\end{figure*}\end{wrapfigure}}
  }

% pdf enhancements
\usepackage[final,colorlinks]{hyperref}[2003/11/30]
% we check if the driver is the useless (for pdf) "hypertex",
% and if so we force pdftex instead and issue a warning
\def\@tempa{hypertex}
\ifx\@tempa\Hy@defaultdriver
  \GenericWarning{*** }{***\MessageBreak
   Inappropriate driver for hyperref specified: assuming pdftex.\MessageBreak
   You should amend the source code if using another driver.\MessageBreak
   \expandafter\@gobble\@gobble}
  \Hy@SetCatcodes\input{hpdftex.def}\Hy@RestoreCatcodes
\fi
\usepackage{memhfixc}[2004/05/13]
\providecommand\ebook{1wxyz} % temp value; will be overwritten by WWer
\hypersetup{pdftitle=The Project Gutenberg eBook \#\ebook: How to Draw a Straight Line,
  pdfsubject=A Lecture on Linkages,
  pdfauthor=A.B. Kempe,
  pdfkeywords={Joshua Hutchinson, David Wilson,
               Project Gutenberg Online Distributed Proofreading Team,
               Cornell University, Nature Series},
  pdfstartview=Fit,
  pdfstartpage=1,
  pdfpagemode=UseNone,
  pdfdisplaydoctitle,
  bookmarksopen,
  bookmarksopenlevel=1,
  linktocpage=false,
  pdfpagescrop=0 0 422 652, % 149x230mm
  pdfpagelayout=TwoPageRight, % this is Acrobat 6's "Facing"
  plainpages=false, linkcolor=\ifdraftdoc blue\else black\fi,
  menucolor=\ifdraftdoc blue\else black\fi,
  citecolor=\ifdraftdoc blue\else black\fi,
  urlcolor=\ifdraftdoc magenta\else black\fi}

% for adding explicit destinations
\newcommand\DPanchor[1]{\rlap{\hyper@anchorstart{#1}\hyper@anchorend}}
% to add an explicit level 0 bokmark
\def\ForceBookmark#1#2{\llap{\Hy@writebookmark{}{#1}{#2}{0}{toc}}}

% A quasi-verbatim environment for boilerplate, slightly less drastic than alltt
% Spaces, linebreaks, $ , % and # will appear as typed
% but unlike full verbatim, commands will still be interpreted and long lines will wrap
% (comments documenting the boilerplate text need to use | as the comment character)
% Uses slightly non-standard obeylines and active space.
% The optional argument can be used to specify an explicit font size for the boilerplate.
% If no optional argument is provided, a fontsize will be computed to allow--as nearly as
% possible--73 fixed-width characters per \textwidth (the longest line in the PG license
% had 73 characters at one stage)
{\catcode`\^^M=\active % these lines must end with %
  \global\def\PGobeylines{\catcode`\^^M\active \def^^M{\null\par}}}%
{\obeyspaces%
\global\def\PGb@ilerplate[#1]{\def\PGb@ilerplateHook{#1}\catcode`\%11\relax%
\catcode`\$11\relax\catcode`\#11\relax\catcode`\|=14\relax%
\pretolerance=\@m\hyphenpenalty=5000%
\rightskip=\z@\@plus20em\relax%
\frenchspacing\ttfamily\PGb@ilerplateHook%
\def {\noindent\null\space}%
\parindent=\z@\PGobeylines\obeyspaces}}
\def\PGboilerplate{%
 \@ifnextchar[{\PGb@ilerplate}{\PGb@ilerplate[\PGAutoFit{73}]}}
\let\endPGboilerplate\empty
% \PGAutoFit adjusts the fontsize so a specified number of
% fixed-width characters will fit in the current \textwidth
\def\PGrem@pt#1.#2Q@!!@Q{#1}
\def\PGAutoFit#1{\setbox\z@=\hbox{m}\dimen@=\wd\z@\relax
  \multiply\dimen@#1\relax \dimen@i=\dimen@\relax
  \dimen@=\textwidth\relax
  \dimen@ii=\f@size pt  \advance\dimen@ii0.5pt
  \expandafter\multiply\expandafter\dimen@\expandafter\PGrem@pt\the\dimen@ii Q@!!@Q
  \expandafter\divide\expandafter\dimen@\expandafter\PGrem@pt\the\dimen@i Q@!!@Q
  \dimen@i=\dimen@\multiply\dimen@i12\divide\dimen@i10
  \fontsize{\strip@pt\dimen@}{\strip@pt\dimen@i}\ttfamily\selectfont}
{\catcode`\L\active
\gdef\PGlicencelink{\catcode`\L\active\letL\PGlinklicence}}
\def\PGlinklicence{\@ifnextchar i{\PG@lli}{L}}
\def\PG@lli#1{\@ifnextchar c{\PG@llii}{Li}}
\def\PG@llii#1{\@ifnextchar e{\PG@lliii}{Lic}}
\def\PG@lliii#1{\@ifnextchar n{\PG@lliv}{Lice}}
\def\PG@lliv#1{\@ifnextchar s{\PG@llv}{Licen}}
\def\PG@llv#1{\@ifnextchar e{\PG@llvi}{Licens}}
\def\PG@llvi#1{\hyperlink{PGlicence}{License}}

\makeatother

% to handle the endnotes in the text, with automatic hyperlinking back and forth
\newcommand\Note[1]{~(\DPanchor{noteb:#1}\hyperlink{note:#1}{#1})}
\newcommand\NoteText[1]{\par\medskip\indent(\DPanchor{note:#1}\hyperlink{noteb:#1}{#1})~\ignorespaces}

% to deal with the scanned page breaks;
% add a "draft" option to the documentclass invocation to see the scan numbers
\ifdraftdoc
\def\PG#1 #2.png#3
{\marginpar{\noindent\null\hfill\Small #2.png}}
\def\PGx#1 #2.png#3
{}
\else
\def\PG#1 #2.png#3
{}
\let\PGx\PG
\fi

% other minor tweaks
\lefthyphenmin=3
\sloppy

\begin{document}
\frontmatter
\pagestyle{empty}

\makehalftitlepage

\iffalse
\PG--File: 001.png---\********\*******\********\********\*********\--------
% Production note
PRODUCTION NOTE

Cornell University Library
produced this volume to replace
the irreparably deteriorated
original. It was scanned using
Xerox software and equipment at
600 dots per inch resolution
and compressed prior to storage
using CCITT Group 4
compression. The digital data
were used to create Cornell's
replacement volume on paper
that meets the ANSI Standard
Z39.48-1984. The production of
this volume was supported in
part by the Commission on
Preservation and Access and the
Xerox Corporation. 1990.

\PG--File: 002.png---\********\*******\********\********\*********\--------
% [Blank Page]
\PG--File: 003.png---\********\*******\********\********\*********\--------
% Title page

% [Illustration: NATURE SERIES]

HOW TO DRAW A STRAIGHT LINE;

A

LECTURE ON LINKAGES.

\PG--File: 004.png---\********\*******\********\********\*********\--------
% [Illustration]
\PG--File: 005.png---\********\******\********\********\*********\---------
% Title page

<i>NATURE SERIES</i>.

HOW TO DRAW A STRAIGHT LINE;

A

LECTURE ON LINKAGES.

BY

A. B. KEMPE, B.A.,

OF THE INNER TEMPLE, ESQ.;
MEMBER OF THE COUNCIL OF THE LONDON MATHEMATICAL SOCIETY;
AND LATE SCHOLAR OF TRINITY COLLEGE, CAMBRIDGE.

WITH NUMEROUS ILLUSTRATIONS.

London:

MACMILLAN AND CO,

1877.

[<i>The Right of Translation and Reproduction is Reserved.</i>]
\fi
% NATURE SERIES in preamble
\title{HOW TO DRAW A STRAIGHT LINE\Huge;\\[28pt]
\normalfont\tiny
\protect\hypertarget{titlepage}{A}\protect\ForceBookmark{Title Page}{titlepage}\\[28pt]
\large
LECTURE ON LINKAGES.}
\author{A. B. KEMPE, B.A.,}
% BY in preamble
\affiliation{OF THE INNER TEMPLE, ESQ.;\\
MEMBER OF THE COUNCIL OF THE LONDON MATHEMATICAL SOCIETY;\\
AND LATE SCHOLAR OF TRINITY COLLEGE, CAMBRIDGE.}
\subtitle{WITH NUMEROUS ILLUSTRATIONS.}
\date{\protect\includegraphics
  [height=11pt,viewport=0 20 182 72]{images/london.pdf}\;: \\ % graphic draws outside boundingbox
  1877.}

\maketitle

\clearpage

\transcribersnote{The original book was published by MacMillan and Co.,
 and printed by R.~Clay, Sons, and Taylor, Printers, Bread Street Hill,
 Queen Victoria Street.}

\transcribersnote{Inconsistent spelling (\hyperlink{qip}{quadriplane}/\hyperlink{qap}{quadruplane}) has been retained.
Some illustrations have been removed slightly from their original locations to avoid infelicitous page breaks.
Minor typographical corrections are documented in the \LaTeX\ source.}
\makecopyrightpage

\iffalse
\PG--File: 006.png---\********\******\********\********\*********\---------
% Title page verso

LONDON:

R. CLAY, SONS, AND TAYLOR, PRINTERS,

BREAD STREET HILL,

QUEEN VICTORIA STREET.
\PG--File: 007.png---\********\******\********\********\*********\---------
\fi
\chapter*{}
\subsection*{\protect\hypertarget{notice}{NOTICE}.\protect\ForceBookmark{Notice}{notice}}

\textsc{This} Lecture was one of the series delivered to science
teachers last summer in connection with the Loan
Collection of Scientific Apparatus. I have taken the
opportunity afforded by its publication to slightly
enlarge it and to add several notes. For the illustrations
I am indebted to my brother, Mr.\ \textsc{H.~R. Kempe},
without whose able and indefatigable co-operation in
drawing them and in constructing the models furnished
by me to the Loan Collection I could hardly have
undertaken the delivery of the Lecture, and still less
its publication.

\medskip
 7, \textsc{Crown Office Row, Temple,}

\qquad \textit{January $16$th}, 1877.
\iffalse
\PG--File: 008.png---\********\******\********\********\*********\---------
% [Blank Page]
\PG--File: 009.png---\********\******\********\********\*********\---------
% half title page
HOW TO DRAW A STRAIGHT LINE;

A

LECTURE ON LINKAGES.
\PG--File: 010.png---\********\******\********\********\*********\---------
% [Blank Page]
\PG--File: 011.png---\********\******\********\********\*********\---------
\fi

\mainmatter
\pagestyle{kempe}

\chapter*{\protect\hypertarget{main}{HOW TO DRAW A STRAIGHT LINE}\Huge:\protect\ForceBookmark{Text}{main}}

\section*{A LECTURE ON LINKAGES.}

\noindent\textsc{The} great geometrician Euclid, before demonstrating to
us the various propositions contained in his \textit{Elements of
Geometry}, requires that we should be able to effect certain
processes. These \emph{Postulates}, as the requirements are
termed, may roughly be said to demand that we should
be able to describe straight lines and circles. And
so great is the veneration that is paid to this master-geometrician,
that there are many who would refuse the
designation of ``geometrical'' to a demonstration which
requires any other construction than can be effected by
straight lines and circles. Hence many problems---such as,
for example, the trisection of an angle---which can readily
be effected by employing other simple means, are said to
have no geometrical solution, since they cannot be solved
by straight lines and circles only.

It becomes then interesting to inquire how we can effect
these preliminary requirements, how we can describe these
circles and these straight lines, with as much accuracy
as the physical circumstances of the problems will admit of.
\PG--File: 012.png---\********\******\********\********\*********\---------

As regards the circle we encounter no difficulty. Taking
Euclid's definition, and assuming, as of course we must,
that our surface on which we wish to describe the circle is
a plane,\Note{1}\footnote
  {These figures refer to Notes at the end of the lecture.} we see that we have only to make our tracing-point
preserve a distance from the given centre of the circle
constant and equal to the required radius. This can readily
be effected by taking a flat piece of any form, such as the
piece of cardboard I have here, and passing a pivot which
is fixed to the given surface at the given centre through a
hole in the piece, and a tracer or pencil through another
hole in it whose distance from the first is equal to the
given radius; we shall then, by moving the pencil, be able,
even with this rude apparatus, to describe a circle with considerable
accuracy and ease; and when we come to employ
very small holes and pivots, or even large ones, turned with
all that marvellous truth which the lathe affords, we shall
get a result unequalled perhaps among mechanical apparatus
for the smoothness and accuracy of its movement. The
apparatus I have just described is of course nothing but a
simple form of a pair of compasses, and it is usual to say
that the third Postulate postulates the compasses.

But the straight line, how are we going to describe that?
Euclid defines it as ``lying evenly between its extreme
points.'' This does not help us much. Our text-books say
that the first and second Postulates postulate a ruler\Note{2}.
But surely that is begging the question. If we are to draw
a straight line with a ruler, the ruler must itself have a
straight edge; and how are we going to make the edge
straight? We come back to our starting-point.

Now I wish you clearly to understand the difference
\PG--File: 013.png---\********\******\********\********\*********\---------
between the method I just now employed for describing a
circle, and the ruler method of describing a straight line.
If I applied the ruler method to the description of a circle,
I should take a circular lamina, such as a penny, and trace
my circle by passing the pencil round the edge, and I should
have the same difficulty that I had with the straight-edge,
for I should first have to make the lamina itself circular.
But the other method I employed involves no begging the
question. I do not first assume that I have a circle and
then use it to trace one, but simply require that the distance
between two points shall be invariable. I am of course
aware that we do employ circles in our simple compass, the
pivot and the hole in the moving piece which it fits are
such; but they are used not because they are the curves
we want to describe (they are not so, but are of a different
size), as is the case with the straight-edge, but because,
through the impossibility of constructing pivots or holes of
no finite dimensions, we are forced to adopt the best substitute
we can for making one point in the moving piece
remain at the same spot. If we employ a very small pivot
and hole, though they be not truly circular, the error in
the description of a circle of moderate dimensions will be
practically infinitesimal, not perhaps varying beyond the
width of the thinnest line which the tracer can be made to
describe; and even when we employ large pivots and holes
we shall get results as accurate, because those pivots and
holes may be made by the employment of very small ones
in the machine which makes them.

It appears then, that although we have an easy and accurate
method of describing a circle, we have at first sight no
corresponding means of describing a straight line; and
\PG--File: 014.png---\********\******\********\********\*********\---------
there would seem to be a substantial difficulty in producing
what mathematicians call the simplest curve, so that the
question how to get over that difficulty becomes one of a
decided theoretical interest.

Nor is the interest theoretical only, for the question is
one of direct importance to the practical mechanician. In
a large number of machines and scientific apparatus it is
requisite that some point or points should move accurately
in a straight line with as little friction as possible. If the
ruler principle is adopted, and the point is kept in its path
by guides, we have, besides the initial difficulty of making
the guides truly straight, the wear and tear produced by
the friction of the sliding surfaces, and the deformation
produced by changes of temperature and varying strains.
It becomes therefore of real consequence to obtain, if
possible, some method which shall not involve these objectionable
features, but possess the accuracy and ease
of movement which characterise our circle-producing
apparatus.

Turning to that apparatus, we notice that all that is
requisite to draw with accuracy a circle of any given
radius is to have the distance between the pivot and the
tracer properly determined, and if I pivot a second ``piece''
to the fixed surface at a second point having a tracer as
the first piece has, by properly determining the distance
between the second tracer and pivot, I can describe a second
circle whose radius bears any proportion I please to that of
the first circle. Now, removing the tracers, let me pivot
a third piece to these two \textit{radial} pieces, as I may call them,
at the points where the tracers were, and let me fix a tracer
at any point on this third or \textit{traversing} piece. You will
\PG--File: 015.png---\******\******\********\********\*********\-----------
at once see that if the radial pieces were big enough the
tracer would describe circles or portions of circles on \textit{them},
though they are in motion, with the same ease and accuracy
as in the case of the simple circle-drawing apparatus; the
tracer will not however describe a circle on the \textit{fixed} surface,
but a complicated curve.

\begin{figure*}[hbt]
\Figure{1}{015}
% [Illustration: Fig. 1.]
\end{figure*}

This curve will, however, be described with all the ease
and accuracy of movement with which the circles were described,
and if I wish to reproduce in a second apparatus the
curves which I produce with this, I have only to get the
distances between the pivots and tracers accurately the same
in both cases, and the curves will also be accurately the same.
I could of course go on adding fresh pieces \textit{ad libitum}, and
I should get points on the structure produced, describing in
general very complicated curves, but with the same results
as to accuracy and smoothness, \textit{the reproduction of any particular
curve depending solely on the correct determination of a
certain definite number of distances}.

These systems, built up of pieces pointed or pivoted
together, and turning about pivots attached to a fixed base,
so that the various points on the pieces all describe definite
\PG--File: 016.png---\******\******\********\********\*********\-----------
curves, I shall term ``link-motions,'' the pieces being termed
``links.'' As, however, it sometimes facilitates the consideration
of the properties of these structures to regard
them apart from the base to which they are pivoted, the
word "linkage" is employed to denote any combination of
pieces pivoted together. When such a combination is
pivoted in any way to a fixed base, the motion of points on
it not being necessarily confined to fixed paths, the link-structure
is called a ``linkwork:'' a ``linkwork'' in which
the motion of every point is in some definite path being,
as before stated, termed a ``link-motion.'' I shall only add
to these expressions two more: the point of a link-motion
which describes any curve is called a ``graph,'' the curve
being called a ``gram''\Note{3}.

The consideration of the various properties of these ``linkages''
has occupied much attention of late years among
mathematicians, and is a subject of much complexity and
difficulty. With the purely mathematical side of the
question I do not, however, propose to deal to-day, as we
shall have quite enough to do if we confine our attention to
the practical results which mathematicians have obtained,
and which I believe only mathematicians could have obtained.
That these results are valuable cannot I think
be doubted, though it may well be that their great beauty
has led some to attribute to them an importance which they
do not really possess; and it may be that fifty years ago
they would have had a value which, through the great improvements
that modern mechanicians have effected in the
production of true planes, rulers and other exact mechanical
structures, cannot now be ascribed to them. But linkages
have not at present, I think, been sufficiently put before
\PG--File: 017.png---\******\*******\********\********\*********\----------
the mechanician to enable us to say what value should
really be set upon them.

The practical results obtained by the use of linkages
are but few in number, and are closely connected with the
problem of ``straight-line motion,'' having in fact been
discovered during the investigation of that problem, and I
shall be naturally led to consider them if I make ``straight-line
motion'' the backbone of my lecture. Before, however,
plunging into the midst of these linkages it will be useful
to know how we can practically construct such models as
we require; and here is one of the great advantages of our
subject---we can get our results visibly before us so very
easily. Pins for fixed pivots, cards for links, string or cotton
for the other pivots, and a dining-room table, or a drawing-board
if the former be thought objectionable, for a fixed
base, are all we require. If something more artistic be
preferred, the plan adopted in the models exhibited by me
in the Loan Collection can be employed. The models were
constructed by my brother, Mr. H.~R. Kempe, in the following
way. The bases are thin deal boards painted black;
the links are neatly shaped out of thick cardboard (it is
hard work making them, you have to sharpen your knife
about every ten minutes, as the cardboard turns the edge
very rapidly); the pivots are little rivets made of catgut,
the heads being formed by pressing the face of a heated
steel chisel on the ends of the gut after it is passed through
the holes in the links; this gives a very firm and smoothly-working
joint. More durable links may be made of tin-plate;
the pivot-holes must in this case be punched, and
the eyelets used by bootmakers for laced boots employed
as pivots; you can get the proper tools at a trifling expense
at any large tool shop.
\PG--File: 018.png---\******\******\********\********\*********\-----------

Now, as I have said, the curves described by the
various points on these link-motions are in general very
complex. But they are not necessarily so. By properly
choosing the distances at our disposal we can
make them very simple. But can we go to the fullest
extent of simplicity and get a point on one of them moving
accurately in a straight line? That is what we are going
to investigate.

To solve the problem with our single link is clearly
impossible: all the points on it describe circles. We must
therefore go to the next simple case---our three-link motion.
In this case you will see that we have at our disposal the
distance between the fixed pivots, the distances between the
pivots on the radial links, the distance between the pivots
on the traversing link, and the distances of the tracer from
those pivots; in all six different distances. Can we choose
those distances so that our tracing-point shall move in a
straight line?

The first person who investigated this was that great man
James Watt. ``Watt's Parallel Motion''\Note{4}, invented in
1784, is well known to every engineer, and is employed in
nearly every beam-engine. The apparatus, reduced to its
simplest form, is shown in \figref{2}.

\begin{figure*}[hbt]
\Figure{2}{019}
% [Illustration: Fig.~2.]
\end{figure*}

The radial bars are of equal length,---I employ the word
``length'' for brevity, to denote the distance between the
pivots; the links of course may be of any length or shape,---and
the distance between the pivots or the traversing link
is such that when the radial bars are parallel the line
joining those pivots is perpendicular to the radial bars.
The tracing-point is situate half-way between the pivots
on the traversing piece. The curve described by the tracer
is, if the apparatus does not deviate much from its mean
\PG--File: 019.png---\******\******\********\******\*********\-------------
position, approximately a straight line. The reason of this
is that the circles described by the extremities of the
radial bars have their concavities turned in opposite directions,
and the tracer being half-way between, describes a
curve which is concave neither one way nor the other, and
is therefore a straight line. The curve is not, however,
accurately straight, for if I allow the tracer to describe the
whole path it is capable of describing, it will, when it
gets some distance from its mean position, deviate considerably
from the straight line, and will be found to describe a
figure 8, the portions at the crossing being nearly straight.
We know that they are not quite straight, because it is
impossible to have such a curve partly straight and partly
curved.

For many purposes the straight line described by Watt's
apparatus is sufficiently accurate, but if we require an exact
one it will, of course, not do, and we must try again. Now
it is capable of proof that it is impossible to solve the
problem with three moving links; closer approximations
\PG--File: 020.png---\******\******\********\******\*********\-------------
to the truth than that given by Watt can be obtained, but
still not actual truth.

I have here some examples of these closer approximations.
The first of these, shown in \figref{3}, is due to Richard
Roberts of Manchester.

\begin{figure*}[hbt]
\Figure{3}{020}
% [Illustration: Fig.~3.]
\end{figure*}

The radial bars are of equal length, the distance
between the fixed pivots is twice that of the pivots on the
traversing piece, and the tracer is situate on the traversing
piece, at a distance from the pivots on it equal to the lengths
of the radial bars. The tracer in consequence coincides
with the straight line joining the fixed pivots at those
pivots and half-way between them. It does not, however,
coincide at any other points, but deviates very slightly
between the fixed pivots. The path described by the
tracer when it passes the pivots altogether deviates from
the straight line.
\PG--File: 021.png---\******\**************\********\********\*********\---

The other apparatus was invented by Professor Tchebicheff
of St.\ Petersburg. It is shown in \figref{4}. The radial
\begin{figure*}[hbt]
\Figure{4}{021}
% [Illustration: Fig.~4.]
\end{figure*}
bars are equal in length, being each in my little model
five inches long. The distance between the fixed pivots
must then be four inches, and the distance between the
pivots or the traversing bar two inches. The tracer is
taken half-way between these last. If now we draw a
straight line---I had forgotten that we cannot do that
yet, well, if we draw a straight line, popularly so called---through
the tracer in its mean position, as shown in
the figure, parallel to that forming the fixed pivots, it
will be found that the tracer will coincide with that line at
the points where verticals through the fixed pivots cut it as
well as at the mean position, but, as in the case of Roberts's
parallel motion, it coincides nowhere else, though its deviation
is very small as long as it remains between the verticals.
\PG--File: 022.png---\******\**************\********\********\*********\---

We have failed then with three links, and we must go on
to the next case, a five-link motion---for you will observe
that we must have an odd number of links if we want an
apparatus describing definite curves. Can we solve the
problem with five? Well, we can; but this was not the
first accurate parallel motion discovered, and we must give
the first inventor his due (although he did not find the
simplest way) and proceed in strict chronological order.

In 1864, eighty years after Watt's discovery, the problem
was first solved by M.~Peaucellier, an officer of Engineers
\begin{figure*}[hbt]
\Figure{5}{022}
% [Illustration: Fig.~5.]
\end{figure*}
in the French army. His discovery was not at first estimated
at its true value, fell almost into oblivion, and was
rediscovered by a Russian student named Lipkin, who got
a substantial reward from the Russian Government for his
supposed originality. However, M.~Peaucellier's merit
has at last been recognized, and he has been awarded the
great mechanical prize of the Institute of France, the
``Prix Montyon.''

M.~Peaucellier's apparatus is shown in \figref{5}. It has,
as you see, seven pieces or links. There are first of all
\PG--File: 023.png---\******\**************\********\********\*********\---
two long links of equal length. These are both pivoted at
the same fixed point; their other extremities are pivoted to
opposite angles of a rhombus composed of four equal shorter
links. The portion of the apparatus I have thus far
described, considered apart from the fixed base, is a linkage
termed a ``Peaucellier cell.'' We then take an \emph{extra} link,
and pivot it to a fixed point whose distance from the first
fixed point, that to which the cell is pivoted, is the same as
the length of the extra link; the other end of the extra
link is then pivoted to one of the free angles of the
rhombus; the other free angle of the rhombus has a pencil
at its pivot. That pencil will accurately describe a straight
line.

I must now indulge in a little simple geometry. It is
absolutely necessary that I should do so in order that you
may understand the principle of our apparatus.

In \figref{6}, $Q C$ is the extra link pivoted to the fixed point
\begin{figure*}[htb]
\Figure{6}{024}
% [Illustration: Fig.~6.]
\end{figure*}
$Q$, the other pivot on it $C$, describing the circle
$O C R$. The
straight lines $P M$ and $P' M'$ are supposed to be perpendicular
to $M R Q O M'$.

Now the angle $O C R$, being the angle in a semicircle, is
a right angle. Therefore the triangles $O C R$, $O M P$ are
similar. Therefore,
\[
O C : O R :: O M : O P.
\]
Therefore,
\[
OC\cdot OP = OM\cdot OR,
\]
wherever $C$ may be on the circle. That is, since $O M$ and
$O R$ are both constant, if while $C$ moves in a circle $P$ moves
\PG--File: 024.png---\******\**************\********\******\*********\-----
so that $O$, $C$, $P$ are always in the same straight line, and
so that $OC\cdot OP$ is always constant; then $P$ will describe
the straight line $P~M$ perpendicular to the line $O~Q$.

It is also clear that if we take the point $P'$ on the other
side of $O$, and if $OC\cdot OP'$ is constant $P'$ will describe the
straight line $P'M'$. This will be seen presently to be important.

Now, turning to \figref{7}, which is a skeleton drawing of
\begin{figure*}[hbt]
\Figure{7}{025}
% [Illustration: Fig. 7.]
\end{figure*}
the Peaucellier cell, we see that from the symmetry of the
construction of the cell, $O$, $C$, $P$, all lie in the same straight % "t" not printed in scan
line, and if the straight line $A~n$ be drawn perpendicular
to $C~P$---it must still be an imaginary one, as we have not
proved yet that our apparatus does draw a straight line---$Cn$
is equal to $nP$.
\PG--File: 025.png---\******\**************\*******\********\********\-----

Now,
\begin{align*}
	O~A^2& = On^2 + An^2\\
	A~P^2& = Pn^2 + An^2
\end{align*}
therefore,
\begin{align*}
	O~A^2-A~P^2& = On^2-Pn^2\\
	& = [On-Pn] \cdot [ On + Pn ]\\
	& = OC \cdot OP.
\end{align*}
Thus since $O~A$ and $A~P$ are both constant $OC \cdot OP$ is
always constant, however far or near $C$ and $P$ may be to $O$.
If then the pivot $O$ be fixed to the point $O$ in \figref{6}, and
the pivot $C$ be made to describe the circle in the figure by
being pivoted to the end of the extra link, the pivot $P$ will
satisfy all the conditions necessary to make it move in a
straight line, and if a pencil be fixed at $P$ it will draw a
straight line. The distance of the line from the fixed
pivots will of course depend on the magnitude of the
quantity $OA^2-OP^2$ which may be varied at pleasure.

I hope you clearly understand the two elements composing
the apparatus, the extra link and the cell, and
the part each plays, as I now wish to describe to you some
modifications of the cell. The extra link will remain the
\PG--File: 026.png---\******\**************\********\******\********\------
same as before, and it is only the cell which will undergo
alteration.

If I take the two linkages in \figref{8}, which are known
\begin{figure*}[hbt]
\Figure{8}{026}
% [Illustration: Fig. 8.]
\end{figure*}
as the ``kite'' and the ``spear-head,'' and place one on the
other so that the long links of the one coincide with those
of the other, and then amalgamate the coincident long links
together, we shall get the original cell of \figref[Figs.]{5} and\figref[]{7}.
If then we keep the angles between the long links, or
that between the short links, the same in the ``kite'' and
``spear-head,'' we see that the height of the ``kite''
multiplied by that of the ``spear-head'' is constant.

Let us now, instead of amalgamating the long links of
the two linkages, amalgamate the short ones. We then get
the linkage of \figref{9}; and if the pivot where the short
\begin{figure*}[hbt]
\Figure{9}{027a}
% [Illustration: Fig. 9.]
\end{figure*}
links meet is fixed, and one of the other free pivots be
made to move in the circle of \figref{6} by the extra link,
the other will describe, not the straight line $P~M$, but the
straight line $P'~M'$. In this form, which is a very compact
\PG--File: 027.png---\******\**************\********\******\********\------
one, the motion has been applied in a beautiful manner
to the air engines which are employed to ventilate the
Houses of Parliament. The ease of working and absence
of friction and noise is very remarkable. The engines
were constructed and the Peaucellier apparatus adapted to
them by Mr.\ Prim, the engineer to the Houses, by whose
courtesy I have been enabled to see them, and I can assure
you that they are well worth a visit.

Another modification of the cell is shown in \figref{10}.
If instead of employing a ``kite'' and ``spear-head'' of the
\begin{figure*}[hbt]
\Figure{10}{027b}
% [Illustration: Fig. 10.]
\end{figure*}
same dimensions, I take the same ``kite'' as before, but
use a ``spear-head'' of half the size of the former one,
\PG--File: 028.png---\******\**************\********\******\********\------
the angles being however kept the same, the product of the
heights of the two figures will be half what it was before,
but still constant. Now instead of superimposing the
links of one figure on the other, it will be seen that in % Fig.\
\figref{10} I fasten the shorter links of each figure together, end
to end. Then, as in the former cases, if I fix the pivot at
the point where the links are fixed together, I get a cell
which may be used, by the employment of an extra link, to
describe a straight line. A model employing this form of
cell is exhibited in the Loan Collection by the Conservatoire
des Arts et M�tiers of Paris, and is of exquisite workmanship;
the pencil seems to swim along the straight line.

M.~Peaucellier's discovery was introduced into England
by Professor Sylvester in a lecture he delivered at the
Royal Institution in January, 1874\Note{5}, which excited very
great interest and was the commencement of the consideration
of the subject of linkages in this country.

In August of the same year Mr.\ Hart of Woolwich
Academy read a paper at the British Association meeting\Note{6},
in which he showed that M.~Peaucellier's cell could be
replaced by an apparatus containing only four links instead
of six. The new linkage is arrived at thus.

If to the ordinary Peaucellier cell I add two fresh links
of the same length as the long ones I get the double, or
rather quadruple cell, for it may be used in four different
ways, shown in \figref{11}. Now Mr.\ Hart found that if he
\begin{figure*}[hbt]
\Figure{11}{029a}
% [Illustration: Fig. 11.]
\end{figure*}
took an ordinary parallelogrammatic linkwork, in which the
adjacent sides are unequal, and crossed the links so as to
form what is called a contra-parallelogram, \figref{12}, and
\begin{figure*}[hbt]
\Figure{12}{029b}
% [Illustration: Fig. 12.]
\end{figure*}
then took four points on the four links dividing the distances
between the pivots in the same proportion, those
\PG--File: 029.png---\******\**************\********\******\********\------
four points had exactly the same properties as the four
points of the double cell. That the four points always lie
in a straight line is seen thus: considering the triangle
$abd$, since $aO:Ob:\;:aP:Pd$ therefore $OP$ is parallel to
$bd$, and the perpendicular distance between the parallels is
to the height of the triangle $abd$ as $Ob$ is to $ab$; the same
reasoning applies to the straight line $CO'$, and since $ab:Ob:\;:c~d:O'd$
and the heights of the triangles $abd$, $cbd$, are clearly
the same, therefore the distances of $O~P$ and $O'C$ from $b~d$
are the same, and $O~C~P~O'$ lie in the same straight line.

That the product $OC\cdot OP$ is constant appears at once
\PG--File: 030.png---\******\**************\*******\******\********\-------
when it is seen that $ObC$ is half a ``spear-head'' and $OaP$
half a ``kite;'' similarly it may be shown that $O'P\cdot O'C$ is
constant, as also $OC\cdot CO'$ and $OP\cdot PO'$. Employing then the
Hart's cell as we employed Peaucellier's, we get a five-link
straight line motion. A model of this is exhibited in the
Loan Collection by M.~Breguet.

I now wish to call your attention to an extension of Mr.\
Hart's apparatus, which was discovered simultaneously by
Professor Sylvester and myself. In Mr.\ Hart's apparatus
we were only concerned with bars and points on those bars,
but in the apparatus I wish to bring before you we have
pieces instead of bars. I think it will be more interesting
if I lead up to this apparatus by detailing to you its
history, especially as I shall thereby be enabled to bring
before you another very elegant and very important linkage---the
discovery of Professor Sylvester.

When considering the problem presented by the ordinary
three-\emph{bar} motion consisting of two radial bars and a traversing
bar, it occurred to me---I do not know how or why,
it is often very difficult to go back and find whence one's
ideas originate---to consider the relation between the curves
described by the points on the traversing bar in any given
three-bar motion, and those described by the points on
a similar three-bar motion, but in which the traversing bar
and one of the radial bars had been made to change places.
The proposition was no sooner stated than the solution became
obvious; the curves were precisely similar. In \figref{13}
\begin{figure*}[hbt]
\Figure{13}{031}
% [Illustration: Fig. 13.]
\end{figure*}
let $C~D$ and $B~A$ be the two radial bars turning about the
fixed centres $C$ and $B$, and let $D~A$ be the traversing bar,
and let $P$ be any point on it describing a curve depending
on the lengths of $AB$, $BC$, $CD$, and $DA$. Now add to the
\PG--File: 031.png---\******\**************\********\******\********\------
three-bar motion the bars $CE$ and $EAP'$, $CE$ being equal to
$DA$, and $EA$ equal to $CD$. $C~D~A~E$ is then a parallelogram,
and if an imaginary line $CPP'$ be drawn, cutting $EA$ produced
in $P'$, it will at once be seen that $P'$ is a fixed point
on $EA$ produced, and $CP'$ bears always a fixed proportion to
$CP$, viz., $CD:CE$. Thus the curve described by $P'$ is precisely
the same as that described by $P$, only it is larger in
the proportion $CE:CD$. Thus if we take away the bars
$CD$ and $DA$, we shall get a three-bar linkwork, describing
precisely the same curves, only of different magnitude, as
our first three-bar motion described, and this new three-bar
linkwork is the same as the old with the radial link $CD$
and the traversing link $DA$ interchanged\Note{7}.

On my communicating this result to Professor Sylvester,
he at once saw that the property was one not confined to
\PG--File: 032.png---\******\**************\********\******\********\------
the particular case of points lying on the traversing bar,
in fact to three-\emph{bar} motion, but was possessed by three-\emph{piece}
motion. In \figref{14} $C~D~A~B$ is a three-bar motion, as in
\figref{13}, but the tracing point or ``graph'' does not lie on
the line joining the joints $A~D$, but is anywhere else on a
\begin{figure*}[hbt]
\Figure{14}{032}
% [Illustration: Fig. 14.]
\end{figure*}
``piece'' on which the joints $A~D$ lie. Now, as before, add
the bar $C~E$, $C~E$ being equal to $A~D$, and the piece
$A~E~P'$, making $A~E$ equal to $C~D$, and the triangle
$A~E~P'$ similar to the triangle $P~D~A$; so that the angles
$A~E~P'$, $A~D~P$ are equal, and
\[
P'~E : E~A :\;: A~D : D~P.
\]
It follows easily from this---you can work it out for yourselves
without difficulty---that the ratio $P'~C : P~C$ is constant
and the angle $P~C~P'$ is constant; thus the paths of
$P$ and $P'$, or the ``grams'' described by the ``graphs,'' $P$
\PG--File: 033.png---\******\**************\********\******\********\------
and $P'$, are similar, only they are of different sizes, and one
is turned through an angle with respect to the other.

Now you will observe that the two proofs I have given
are quite independent of the bar $A~B$, which only affects
the particular curve described by $P$ and $P'$. If we get rid
of $A~B$, in both cases we shall get in the first figure the
ordinary pantagraph, and in the second a beautiful extension
of it, called by Professor Sylvester, its inventor, the
\textit{Plagiograph} or \textit{Skew Pantagraph}. Like the Pantagraph, it
will enlarge or reduce figures, but it will do more, it will
turn them through any required angle, for by properly
choosing the position of $P$ and $P'$, the ratio of $C~P$ to
$C~P'$ can be made what we please, and also the angle
$P~C~P'$ can be made to have any required value. If the
angle $P~C~P'$ is made equal to 0 or 180�, we get the two
forms of the pantagraph now in common use; if it be
made to assume successively any value which is a sub-multiple
of 360�, we can, by passing the point $P$ each time
over the same pattern make the point $P'$ reproduce it
round the fixed centre $C$ after the fashion of a kaleidoscope.
I think you will see from this that the instrument, which
has, as far as I know, never been practically constructed,
deserves to be put into the hands of the designer. I give
\hyperlink{fig:15}{here} a picture of a little model of a possible form for the
\begin{figure*}[hbt]
\Figure{15}{034}
% [Illustration: Fig. 15.]
\end{figure*}
instrument furnished by me to the Loan Collection by
request of Professor Sylvester\Note{8}.

After this discovery of Professor Sylvester, it occurred to
him and to me simultaneously---our letters announcing our
discovery to each other crossing in the post---that the
principle of the plagiograph might be extended to Mr.\
Hart's contra-parallelogram; and this discovery I shall
\PG--File: 034.png---\******\**************\********\******\********\------
now proceed to explain to you. I shall, however, be
more easily able to do so by approaching it in a different
manner to that in which I did when I discovered it.

If we take the contra-parallelogram of Mr.\ Hart, and
bend the links at the four points which lie on the same
straight line, or \textit{foci} as they are sometimes termed,
\PG--File: 035.png---\******\**************\********\******\********\------
through the same angle, the four points, instead of
lying in the same straight line, will lie at the four angular
points of a parallelogram of constant angles,---two the angle
that the bars are bent through, and the other two their
supplements---and of constant area, so that the product
of two adjacent sides is constant.

In \figref{16} the lettering is preserved as in \figref{12}, so
that the way in which the apparatus is formed may be at
once seen. The holes are taken in the middle of the links,
\begin{figure*}[hbt]
\Figure{16}{035}
% [Illustration: Fig. 16.]
\end{figure*}
and the bending is through a right angle. The four holes
$O~P~O'~C$ lie at the four corners of a right-angled parallelogram,
and the product of any two adjacent sides, as for
example $O~C\cdot O~P$, is constant. It follows that if $O$ be
pivoted to the fixed point $O$ in \figref{6}, and $C$ be pivoted to
the extremity of the extra link, $P$ will describe a straight
line, not $P~M$, but one inclined to $P~M$ at an angle the same
as the bars are bent through, \textit{i.e.}\ a right angle. Thus the
\PG--File: 036.png---\******\**************\********\******\********\------
straight line will be parallel to the line joining the fixed
pivots $O$ and $Q$. This apparatus, which for simplicity I
have described as formed of four straight links which are
afterwards bent, is of course strictly speaking formed of
four plane links, such as those employed in \figref{1}, on which
the various points are taken. This explains the name
given to it by Professor Sylvester, the ``\hypertarget{qap}{Quadruplane}.'' % spelling clear on scan; cf "quadriplane" below
Its properties are not difficult to investigate, and when I
point out to you that in \figref{16}, as in \figref{12}, $Ob$, $bC$ form
half a ``spear-head,'' and $Oa$, $aP$ half a ``kite,'' you will
very soon get to the bottom of it.

I cannot leave this apparatus, in which my name is
associated with that of Professor Sylvester, without expressing
my deep gratitude for the kind interest which he took
in my researches, and my regret that his departure for
America to undertake the post of Professor in the new
Johns Hopkins University has deprived me of one whose
valuable suggestions and encouragement helped me much
in my investigations.

Before leaving the Peaucellier cell and its modifications,
I must point out another important property they possess
besides that of furnishing us with exact rectilinear motion.
We have seen that our simplest linkwork enables us to
describe a circle of any radius, and if we wished to describe
one of ten miles' radius the proper course would be to have
a ten-mile link, but as that would be, to say the least,
cumbrous, it is satisfactory to know that we can effect our
purpose with a much smaller apparatus. When the Peaucellier
cell is mounted for the purpose of describing a
straight line, as I told you, the distance between the fixed
pivots must be the same as the length of the ``extra'' link.
\PG--File: 037.png---\******\**************\********\******\********\------
If this distance be not the same we shall not get straight
lines described by the pencil, but circles. If the difference
be slight the circles described will be of enormous magnitude,
decreasing in size as the difference increases. If the
distance $Q~O$, \figref{6}, be made greater than $Q~C$, the convexity
of the portion of the circle described by the pencil
(for if the circles are large it will of course be only a
portion which is described) will be towards $O$, if less
the concavity. To a mathematician, who knows that the
inverse of a circle is a circle, this will be clear, but it may
not be amiss to give here a short proof of the proposition.

\begin{figure*}[hbt]
\Figure{17}{038}
% [Illustration: Fig. 17.]
\end{figure*}

In \figref{17} let the centres $Q$, $Q'$ of the two circles be at
distances from $O$ proportional to the radii of the circles.
If then $O~D~C~P~S$ be any straight line through $O$, $D~Q$
will be parallel to $P~Q'$, and $C~Q$ to $S~Q'$, and $O~D$ will
bear the same proportion to $O~P$ that $O~Q$ does to $O~Q'$.
Now considering the proof we gave in connection with % Fig.\
\figref{7}, it will be clear that the product $O~D\cdot O~C$ is constant,
and therefore since $O~P$ bears a constant ratio to $O~D$,
$O~P\cdot O~C$ is constant. That is if $O~C\cdot O~P$ is constant and
$C$ describes a circle about $Q$, $P$ will describe one about $Q'$.
Taking then $O$, $C$ and $P$ as the $O$, $C$ and $P$ of the Peaucellier
cell in \figref{7}, we see how $P$ comes to describe a circle.

It is hardly necessary for me to state the importance of
the Peaucellier compass in the mechanical arts for drawing
circles of large radius. Of course the various modifications
of the ``cell'' I have described may all be employed for
the purpose. The models exhibited by the Conservatoire
and M.~Breguet are furnished with sliding pivots for the
purpose of varying the distance between $O$ and $Q$, and
thus getting circles of any radius.
\PG--File: 038.png---\******\**************\********\******\********\------

My attention was first called to these linkworks by the
lecture of Professor Sylvester, to which I have referred.
A passage in that lecture in which it was stated that there
were probably other forms of seven-link parallel motions
\PG--File: 039.png---\******\**************\********\******\********\------
besides M. Peaucellier's, then the only one known, led me
to investigate the subject, and I succeeded in obtaining
some new parallel motions of an entirely different character
to that of M. Peaucellier\Note{9}. I shall bring two of these
to your notice, as the investigation of them will lead us to
consider some other linkworks of importance.

If I take two kites, one twice as big as the other, such
that the long links of each are twice the length of the short
ones, and make one long link of the small kite lie on
a short one of the large, and a short one of the small on
a long one of the large, and then amalgamate the coincident
links, I shall get the linkage shown in \figref{18}.

\begin{figure*}[hbt]
\Figure{18}{039}
% [Illustration: Fig. 18.]
\end{figure*}

The important property of this linkage is that, although
we can by moving the links about, make the points $P$ and
$P'$ approach to or recede from each other, the imaginary
line joining them is always perpendicular to that drawn
through the pivots on the bottom link $L~M$. It follows
that if either of the pivots $P$ or $P'$ be fixed, and the link
$L~M$ be made to move so as always to remain parallel to
a fixed line, the other point will describe a straight line
perpendicular to the fixed line. \figref{19} shows you the
\begin{figure*}[hbt]
\Figure{19}{040}
% [Illustration: Fig.~19.]
\end{figure*}
parallel motion made by fixing $P'$. It is unnecessary for
\PG--File: 040.png---\******\**************\********\******\******\--------
me to point out how the parallelism of $L~M$ is preserved
by adding the link $S~L$, it is obvious from the figure. The
straight line which is described by the point $P$ is perpendicular
to the line joining the two fixed pivots; we can, however,
without increasing the number of links, make a point
on the linkwork describe a straight line inclined to the line
$S~P$ at any angle, or rather we can, by substituting for the
straight link $P~C$ a plane piece, get a number of points on
that piece moving in every direction.

In \figref{20}, for simplicity, only the link $C~P'$ and the
\begin{figure*}[hbt]
\Figure{20}{041}
% [Illustration: Fig.~20.]
\end{figure*}
new piece substituted for the link $P~C$ are shown. The
new piece is circular and has holes pierced in it all at the
same distance---the same as the lengths $P~C$ and $P'~C$---from
$C$. Now we have seen from \figref{19} that $P$ moves in
a vertical straight line, the distance $P~C$ in \figref{20} being
the same as it was in \figref{19}; but from a well-known
property of a circle, if $H$ be any one of the holes pierced in
the piece, the angle $H~P'~P$ is constant, thus the straight
line $H~P'$ is fixed in position, and $H$ moves along it; similarly
all the other holes move along in straight lines passing
through the fixed pivot $P'$, and we get straight line
\PG--File: 041.png---\******\**************\********\******\******\--------
motion distributed in all directions. This species of motion
is called by Professor Sylvester ``tram-motion.'' It is
worth noticing that the motion of the circular disc is the
same as it would have been if the dotted circle on it rolled
inside the large dotted circle; we have, in fact, White's
parallel motion reproduced by linkwork. Of course, if we
only require motion in one direction, we may cut away all
the disc except a portion forming a bent arm containing $C$,
$P$, and the point which moves in the required direction.

The double kite of \figref{18} may be employed to form
some other useful linkworks. It is often necessary to
\begin{figure*}[hbt]
\Figure{21}{042a}
% [Illustration: Fig.~21.]
\end{figure*}
have, not a single point, but a whole piece moving so that
all points on it move in straight lines. I may instance
\PG--File: 042.png---\******\**************\********\******\******\--------
the slide rests in lathes, traversing tables, punches,
drills, drawbridges, etc. The double kite enables us
to produce linkworks having this property. In the linkwork
of \figref{21}, the construction of which will be at once
\begin{figure*}[hbt]
\Figure{22}{042b}
% [Illustration: Fig.~22.]
\end{figure*}
appreciated if you understand the double kite, the horizontal
link moves to and fro as if sliding in a fixed
\PG--File: 043.png---\******\**************\********\******\******\--------
horizontal straight tube.  This form would possibly be
useful as a girder for a drawbridge.
\begin{figure*}[p]
\Figure{24}{044}
% [Illustration: Fig.~24.]
\end{figure*}

{% warning: this section rather fragile
% hardcoded parameters chosen to squish figure 23 onto the
% same page as figure 22 and allow the full-page figure 24
% to follow immediately (otherwise we get a logjam of
% figures on subsequent pages)
\intextsep=0pt % set figure tight inside text
\begin{wrapfigure}[9]{r}{.4\textwidth}
\Figure{23}{043}
% [Illustration: Fig.~23.]
\end{wrapfigure}
In the linkwork of \figref{22}, which is another combination
of two double kites, the vertical link moves so that all its
points move in horizontal straight lines. There is a modification
of this linkwork which will, I think, be found
interesting. In the linkage in \figref{23}, which, if the thin
links are removed, is a skeleton drawing of \figref{22}, % [** original lacks period after "Fig"]
let
the dotted links be taken away and the thin ones be inserted;
we then get a linkage which has the same property
as that in \figref{22}, but it is seen in its new form to be the
ordinary double parallel ruler with three added links.
\figref{24} is a figure of a double parallel ruler made on this
plan with a slight modification. If the bottom ruler be
held horizontal the top moves vertically up and down the
board, having no lateral movement.

}\begin{figure*}[hbt]
\Figure{25}{045}
% [Illustration: Fig.~25.]
\end{figure*}
While I am upon this sort of movement I may point out
an \hyperlink{fig:25}{apparatus} exhibited in the Loan Collection by Professor
\PG--File: 044.png---\******\**************\********\******\******\--------
% [Illustration: Fig.~24.]
\PG--File: 045.png---\******\**************\********\******\******\--------
Tchebicheff, which bears a strong likeness to a complicated
camp-stool, the seat of which has horizontal motion. The
motion is not strictly rectilinear; the apparatus being---as
will be seen by observing that the thin line in the figure
is of invariable length, and a link might therefore be put
where it is---a combination of two of the parallel motions
of Professor Tchebicheff given in \figref{4}, with some links
added to keep the seat parallel with the base. The
variation of the upper plane from a strictly horizontal
\PG--File: 046.png---\*******\**************\********\******\******\-------
movement is therefore double that of the tracer in the
simple parallel motion.

\begin{figure*}[hbt]
\Figure{26}{046}
% [Illustration: Fig.~26.]
\end{figure*}

\figref{26} shows how a similar apparatus of much simpler
construction, employing the Tchebicheff approximate parallel
motion, can be made. The lengths of the links forming
the parallel motion have been given before (\figref{4}). %[**F2: period added.]
The
distance between the pivots on the moving seat is half that
between the fixed pivots, and the length of the remaining
link is one-half that of the radial links.

An \emph{exact} motion of the same description is shown in \figref{27}.
\begin{figure*}[hbt]
\Figure{27}{047}
% [Illustration: Fig[**missing .] 27.]
\end{figure*}
$O$, $C$, $O'$, $P$ are the four \textit{foci} of the \hypertarget{qip}{quadriplane} shown % spelling clear on scan; cf "Quadruplane" above
in the figure in which the links are bent through a right
angle, so that $O~C\cdot O~P$ is constant, and $C~O~P$ a right
angle. The focus $O$ is pivoted to a fixed point, and $C$ is
made by means of the extra link $Q~C$ to move in a circle of
which the radius $Q~C$ is equal to the pivot distance $O$
\PG--File: 047.png---\*******\**************\********\******\******\-------
$P$ consequently moves in a straight line parallel to $O~Q$, the
five moving pieces thus far described constituting the
Sylvester-Kempe parallel motion. To this are added the
moving seat and the remaining link $R~O'$, the pivot distances
of which, $P~R$ and $R~O'$, are equal to $O~Q$. The seat in consequence
always remains parallel to $Q~O$, and as $P$ moves
accurately in a horizontal straight line, every point on
it will do so also. This apparatus might be used with
advantage where a very smoothly-working traversing table
is required.

\begin{figure*}[hbt]
\Figure{28}{048}
% [Illustration: Fig.~28.]
\end{figure*}
\begin{figure*}[hbt]
\Figure{29}{049}
% [Illustration: Fig. 29[**missing .]]
\end{figure*}

I now come to the second of the parallel motions I said I
would show you. If I take a kite and pivot the blunt end to
the fixed base, and make the sharp end move up and down in
a straight line, passing through the fixed pivot, the short
links will rotate about the fixed pivot with equal velocities
\PG--File: 048.png---\*******\**************\********\******\******\-------
in opposite directions; and, conversely, if the links rotate
with equal velocity in opposite directions, the path of the
sharp end will be a straight line, and the same will hold
good if instead of the short links being pivoted to the
same point they are pivoted to different ones.

To find a linkwork which should make two links rotate
with equal velocities in opposite directions was one of the
\PG--File: 049.png---\*******\**************\********\******\******\-------
first problems I set myself to solve. There was no difficulty
in making two links rotate with equal velocities in the same
direction,---the ordinary parallelogrammatic linkwork employed
in locomotive engines, composed of the engine, the two
cranks, and the connecting rod, furnished that; and there
was none in making two links rotate in opposite directions
\PG--File: 050.png---\*******\**************\********\******\******\-------
with \emph{varying} velocity; the contra-parallelogram gave that;
but the required linkwork had to be discovered. After some
trouble I succeeded in obtaining it by a combination of
a large and small contra-parallelogram put together just as
the two kites were in the linkage of \figref{18}. One contra-parallelogram
is made twice as large as the other, and the
long links of each are twice as long as the short\Note{10}.

\begin{figure*}[hbt]
\Figure{30}{050}
% [Illustration: Fig.~30.]
\end{figure*}
\begin{figure*}[hbt]
\Figure{31}{051}
% [Illustration: Fig.~31.]
\end{figure*}

The linkworks in \hyperlink{fig:30}{Figs.\ 30} and \hyperlink{fig:31}{31} will, by considering
the thin line drawn through the fixed pivots in each as a
link, be seen to be formed by fixing different links of the
same six-link linkage composed of two contra-parallelograms
as just stated. The pointed links rotate with equal velocity
in opposite directions, and thus, as shown in \figref{28}, at
once give parallel motions. They can of course, however,
be usefully employed for the mere purpose of reversing
angular velocity\Note{11}.

An extension of the linkage employed in these two last
\PG--File: 051.png---\*******\**************\********\******\******\-------
figures gives us an apparatus of considerable interest. If
I take another linkage contra-parallelogram of half the
size of the smaller one and fit it to the smaller exactly as
I fitted the smaller to the larger, I get the eight-linkage of
\begin{figure*}[hbt]
\Figure{32}{052}
% [Illustration: Fig.~32.]
\end{figure*}
\figref{32}. It has, you see, four pointed links radiating from
a centre at equal angles; if I open out the two extreme
\PG--File: 052.png---\*******\**************\********\******\******\-------
ones to any desired angle, you will see that the two intermediate
ones will exactly \emph{trisect the angle}. Thus the power
we have had to call into operation in order to effect Euclid's
first Postulate---linkages---enables us to solve a problem
\PG--File: 053.png---\*******\**************\********\******\******\-------
\begin{figure*}[hbt]
\Figure{33}{053}
% [Illustration: Fig.~33[**missing .]]
\end{figure*}
which has no ``geometrical'' solution. I could of course
go on extending my linkage and get others which would
divide an angle into any number of equal parts. It is obvious
\PG--File: 054.png---\*******\**************\********\******\******\-------
that these same linkages can also be employed as linkworks
for doubling, trebling, etc., angular velocity\Note{12}.

Another form of ``Isoklinostat''---for so the apparatus is
termed by Professor Sylvester---was discovered by him. The
construction is apparent from \figref{33}. It has the great
advantage of being composed of links having only two
pivot distances bearing any proportion to each other, but
it has a larger number of links than the other, and as the
opening out of the links is limited, it cannot be employed
for multiplying angular motion.

Subsequently to the publication of the paper which contained
an account of these linkworks of mine of which I
have been speaking, I pointed out in a paper read before
the Royal Society\Note{13} that the parallel motions given there
were, as well as those of M.~Peaucellier and Mr.~Hart, all
particular cases of linkworks of a very general character,
all of which depended on the employment of a linkage
composed of two similar figures. I have not sufficient time,
and I think the subject would not be sufficiently inviting
on account of its mathematical character, to dwell on it
here, so I will leave those in whom an interest in the
question has been excited to consider the original paper.

At this point the problem of the production of straight-line
motion now stands, and I think you will be of opinion
that we hardly, for practical purposes, want to go much
farther into the theoretical part of the question. The
results that have been obtained must now be left to the
mechanician to deal with, if they are of any practical value.
I have, as far as what I have undertaken to bring before
you to-day is concerned, come to the end of my tether.
I have shown you that we \emph{can} describe a straight line, and
\PG--File: 055.png---\*******\**************\********\******\******\-------
\emph{how} we can, and the consideration of the problem has led us
to investigate some important pieces of apparatus. But I
hope that this is not all. I hope that I have shown you (and
your attention makes that hope a belief) that this new field
of investigation is one possessing great interest and importance.
Mathematicians have no doubt done much more
than I have been able to show you to-day\Note{14}, but the unexplored
fields are still vast, and the earnest investigator
can hardly fail to make new discoveries. I hope therefore
that you whose duty it is to extend the domain of science
will not let the subject drop with the close of my lecture.
\PG--File: 056.png---\*******\**************\********\******\******\-------
% [Blank Page]
\PG--File: 057.png---\*******\**************\********\******\******\-------

\chapter*{}
\subsection*{\protect\hypertarget{notes}{NOTES}.\protect\ForceBookmark{Notes}{notes}}

{\small
\NoteText{1} The hole through which the pencil passes can be made to describe
a circle independently of any surface (see the latter part of Note~3),
but when we wish to describe a circle or a given plane surface that
surface must of course be assumed to be plane.

\NoteText{2} ``But'' (it is carefully added) ``not a graduated one.'' By the
use of a ruler with only two graduations, an angle can, as is well known,
be readily trisected, thus---Let $RST$ be the angle, and let $PP'$ be the
points where the graduations cut the edge of the ruler. Let $2RS=PP'$.
Draw $RU$ parallel and $RV$ perpendicular to $ST$. Then if we fit the
ruler to the figure $RSTUV$ so that the edge $PP'$ passes through $S$,
$P$ lies on $RU$ and $P'$ on $RV$, $PP'$ trisects the angle $RST$. For if $Q$ be
the middle point of $PP'$, and $RQ$ be joined, the angle $TSP=$ the angle
$QPR=$ the angle $QRP=$ half the angle $RQS$, that is half the angle
$RSQ$.

This solution is of course not a ``geometrical'' one in the sense I
have indicated, because a graduated ruler and the fitting process are
employed. But does Euclid confine himself to his three Postulates of
construction? Does he not use a graduated ruler and this fitting process?
Is not the side $AB$ of the triangle $ABC$ in Book~I. Proposition
4, graduated at $A$ and $B$, and are we not told to take it up and fit it
on to $DE$?

It seems difficult to see why Euclid employed the second Postulate---that
which requires ``that a terminated straight line may be
produced to any length in a straight line,''---or rather, why he did
not put it among the propositions in the First Book as a problem. It is
by no means difficult by a rigid adherence to Euclid's methods to find a
point outside a terminated straight line which is in the same straight
line with it, and to prove it to be so, without the employment of the
second Postulate. That point can then, by the first Postulate, be
\PG--File: 058.png---\********\**************\********\******\******\------
joined to the extremity of the given straight line which is thus produced,
and the process can be continued indefinitely, since by the third
Postulate circles can be drawn with any centre and radius.

\NoteText{3} It is important to notice that the fixed base to which the pivots
are attached is really one link in the system. It would on that account
be perhaps more scientific, in a general consideration of the subject, to
commence by calling any combination of pieces (whether those pieces
be cranks, beams, connecting-rods, or anything else) jointed or pivoted
together, a ``\textit{linkage}.'' When the motion of the links is confined to
one plane or to a number of parallel planes, the system is called
a ``\textit{plane linkage}.'' (It will be seen that this lecture is confined
to plane linkages; a few remarks about solid linkages will be found
at the end of the note.) The motion of the links among themselves
in a linkage may be determinate or not. When the motion is
determinate the number of links must be even, and the linkage is
said to be ``\textit{complete}.'' When the motion is not determinate the
linkage is said to have 1, 2, 3, etc.\ degrees of freedom, according to the
amount of liberty the links possess in their relative motion. These
linkages may be termed ``\textit{primary},'' ``\textit{secondary},'' etc.\ linkages. Thus
if we take the linkage composed of four links with two pivots on
each, the motion of each link as regards the others is determinate, and
the linkage is a ``\textit{complete linkage}.'' If one link be jointed in the
middle the linkage has one degree of liberty and is a ``\textit{primary linkage}.''
So by making fresh joints ``\textit{secondary}'' or ``\textit{tertiary},'' etc.\ linkages
may be formed. These primary, etc.\ linkages may be formed in
various other ways, but the example given will illustrate the reason
for the nomenclature. When one link of a linkage is a fixed base the
structure is called a ``\textit{linkwork}.'' Linkworks, like linkages, may be
``\textit{primary},'' ``\textit{secondary},'' etc. A ``\textit{complete linkwork},'' \textit{i.e.}\ one in
which the motion of every point on the moving part of the structure is
definite, is called a ``\textit{link-motion}.'' The various ``grams'' described by
these link-motions are very difficult to deal with. I have shown, in a
paper in the \textit{Proceedings of the London Mathematical Society}, 1876,
that a link-motion can be found to describe any given algebraic curve,
but the converse problem, ``Given the link-motion, what is the curve?''
is one towards the solution of which but little way has been made; and
the ``tricircular trinodal sextics,'' which are the ``grams'' of the simple
three-piece motion, are still under the consideration of some of our
most eminent mathematicians.

Taking them in their greatest generality, the theoretically simplest
\PG--File: 059.png---\*******\**************\********\******\******\-------
form of link-motion is not the flat circle-producing link, but a solid
link pivoted to a fixed centre, and capable of motion in all directions
about that centre, so that all points on it describe spheres in space;
and the most general form a number of such links pivoted together,
forming a structure the various points on which describe surfaces. If
two simple solid links, turning about two fixed centres, are pivoted
together at a common point, that point will describe a circle independently
of any plane surface, the other points on the links describing
portions of spheres. The form of pivot which would have to be
adopted in solid linkages would be the ball-and-socket joint, so that the
links could not only move about round the fixed centre, but rotate
about any imaginary axis through that centre. It is obvious that it
would be impossible to construct any joint which would give the links
perfect freedom of motion, as the fixed centre about which any link
turned must be fastened to a fixed base in some way, and whatever
means were adopted would interfere with the link in some portion of
its path. This is not so in plane link-motions. The subject of solid
linkages has been but little considered. Hooke's joint may be mentioned
as an example of a solid link-motion. (See also Note~11.)

\NoteText{4} I have been more than once asked to try and get rid of the
objectionable term ``parallel motion.'' I do not know how it came to
be employed, and it certainly does not express what is intended. The
apparatus does not give ``parallel motion,'' but approximate ``rectilinear
motion.'' The expression, however, has now become crystallised,
and I for one cannot undertake to find a solvent.

\NoteText{5} See the \textit{Proceedings of the Royal Institution}, 1874.

\NoteText{6} This paper is printed \textit{in extenso} in the \textit{Cambridge Messenger of
Mathematics}, 1875, vol.~iv., pp.~82--116, and contains much valuable
matter about the mathematical part of the subject.

\NoteText{7} The interchange of a radial and traversing bar converts Watt's
Parallel Motion into the Grasshopper Parallel Motion. The same
change shows us that the curves traced by the linkwork formed by
fixing one bar of a ``kite'' are the same as those traced by the linkwork
formed by fixing one bar of a contra-parallelogram. This is
interesting as showing that there is really only one case in which the
sextic curve, the ``gram'' of three-bar motion, breaks up into a circle
and a quartic.

\NoteText{8} For a full account of this and the piece of apparatus next
described, see \textit{Nature}, vol.~xii., pp.\ 168 and 214.
\PG--File: 060.png---\**********\**************\********\******\******\----

\NoteText{9} See the \textit{Messenger of Mathematics}, ``On Some New Linkages,''
1875, vol.~iv., p.~121.

\NoteText{10} A reference to the paper referred to in the last note will show
that it is not necessary that the small kites and contra-parallelograms
should be half the size of the large ones, or that the long links should
be double the short; the particular lengths are chosen for ease of
description in lecturing.

\NoteText{11} By an arrangement of Hooke's joints, pure solid linkages, we
can make two axes rotate with equal velocities in contrary directions
(See Willis's \textit{Principles of Mechanism}, 2nd~Ed.\ sec.~516, p.~456), and
therefore produce an exact parallel motion.

\NoteText{12} The ``kite'' and the ``contra-parallelogram'' are subject to the
inconvenience (mathematically very important) of having ``dead
points.'' These can be, however, readily got rid of by employing %[**missing g][F2: added]
pins and gabs in the manner pointed out by Professor Reuleaux. (See
Reuleaux's \textit{Kinematics of Machinery}, translated by Professor Kennedy,
Macmillan, pp.~290--294.)

\NoteText{13} ``Proceedings of the Royal Society'', No.~163, 1875, ``On a General
Method of Obtaining Exact Rectilinear Motion by Linkwork.'' I take
this opportunity of pointing out that the results there arrived at may
be greatly extended from the following simple consideration.
% [Illustration: Fig.~34.]

If the straight link $O~B$ makes any angle $D$ with the straight link
$O~A$, and if instead of employing the straight links we employ the pieces
$A'OA$, $B'OB$, and if the angle $A'OA$ equals the angle $B'OB$, then the
\PG--File: 061.png---\**********\**************\********\******\******\----
angle $B'OA'$ equals $D$. The recognition of this very obvious fact will
enable us to derive the Sylvester-Kempe parallel motion from that of
Mr.~Hart.
$$\Figure{34}{060}$$

\NoteText{14} In addition to the authorities already mentioned, the following
may be referred to by those who desire to know more about the mathematical
part of the subject of ``Linkages.'' ``\textit{Sur les Syst�mes de Tiges
Articul�es},'' par M.~V. Liguine, in the \textit{Nouvelles Annales}, December,
1875, pp.~520--560.

Two papers ``\textit{On Three-bar Motion}'' by Professor Cayley and Mr.\
S.~Roberts, in the \textit{Proceedings of the London Mathematical Society},
1876, vol.~vii., pp.~14 and 136. Other short papers in the \textit{Proceedings
of the London Mathematical Society}, vols.~v., vi., vii., and the
\textit{Messenger of Mathematics}, vols.~iv.\ and v. % scan clearly has "Vols" but for consistency we use "vols"

}\iffalse
\PG--File: 062.png---\**********\**************\********\***\******\-------

\begin{center}
{\footnotesize
R. CLAY, SONS, AND TAYLOR, PRINTERS,\\
BREAD STREET HILL, QUEEN VICTORIA STREET.
}
\end{center}

\fi

% we *do* want the licence to start recto, to emphasise it is an addition
\cleartorecto
\pagestyle{licence}
\setlength\parskip{0pt}\raggedbottom
\hypertarget{PGlicence}{ }
\ForceBookmark{Licensing Information}{PGlicence}\par
\begin{PGboilerplate}
End of Project Gutenberg's How to Draw a Straight Line, by A.B. Kempe

*** END OF THIS PROJECT GUTENBERG EBOOK HOW TO DRAW A STRAIGHT LINE ***

***** This file should be named 25155-pdf.pdf or 25155-pdf.zip *****
This and all associated files of various formats will be found in:
        http://www.gutenberg.org/2/5/1/5/25155/

Produced by Joshua Hutchinson, David Wilson and the Online
Distributed Proofreading Team at http://www.pgdp.net (This
file was produced from images from the Cornell University
Library: Historical Mathematics Monographs collection.)


Updated editions will replace the previous one--the old editions
will be renamed.

Creating the works from public domain print editions means that no
one owns a United States copyright in these works, so the Foundation
(and you!) can copy and distribute it in the United States without
permission and without paying copyright royalties.  Special rules,
set forth in the General Terms of Use part of this license, apply to
copying and distributing Project Gutenberg-tm electronic works to
protect the PROJECT GUTENBERG-tm concept and trademark.  Project
Gutenberg is a registered trademark, and may not be used if you
charge for the eBooks, unless you receive specific permission.  If you
do not charge anything for copies of this eBook, complying with the
rules is very easy.  You may use this eBook for nearly any purpose
such as creation of derivative works, reports, performances and
research.  They may be modified and printed and given away--you may do
practically ANYTHING with public domain eBooks.  Redistribution is
subject to the trademark license, especially commercial
redistribution.



*** START: FULL LICENSE ***

THE FULL PROJECT GUTENBERG LICENSE
PLEASE READ THIS BEFORE YOU DISTRIBUTE OR USE THIS WORK

To protect the Project Gutenberg-tm mission of promoting the free
distribution of electronic works, by using or distributing this work
(or any other work associated in any way with the phrase "Project
Gutenberg"), you agree to comply with all the terms of the Full Project
Gutenberg-tm License (available with this file or online at
http://gutenberg.org/license).


Section 1.  General Terms of Use and Redistributing Project Gutenberg-tm
electronic works

1.A.  By reading or using any part of this Project Gutenberg-tm
electronic work, you indicate that you have read, understand, agree to
and accept all the terms of this license and intellectual property
(trademark/copyright) agreement.  If you do not agree to abide by all
the terms of this agreement, you must cease using and return or destroy
all copies of Project Gutenberg-tm electronic works in your possession.
If you paid a fee for obtaining a copy of or access to a Project
Gutenberg-tm electronic work and you do not agree to be bound by the
terms of this agreement, you may obtain a refund from the person or
entity to whom you paid the fee as set forth in paragraph 1.E.8.

1.B.  "Project Gutenberg" is a registered trademark.  It may only be
used on or associated in any way with an electronic work by people who
agree to be bound by the terms of this agreement.  There are a few
things that you can do with most Project Gutenberg-tm electronic works
even without complying with the full terms of this agreement.  See
paragraph 1.C below.  There are a lot of things you can do with Project
Gutenberg-tm electronic works if you follow the terms of this agreement
and help preserve free future access to Project Gutenberg-tm electronic
works.  See paragraph 1.E below.

1.C.  The Project Gutenberg Literary Archive Foundation ("the Foundation"
or PGLAF), owns a compilation copyright in the collection of Project
Gutenberg-tm electronic works.  Nearly all the individual works in the
collection are in the public domain in the United States.  If an
individual work is in the public domain in the United States and you are
located in the United States, we do not claim a right to prevent you from
copying, distributing, performing, displaying or creating derivative
works based on the work as long as all references to Project Gutenberg
are removed.  Of course, we hope that you will support the Project
Gutenberg-tm mission of promoting free access to electronic works by
freely sharing Project Gutenberg-tm works in compliance with the terms of
this agreement for keeping the Project Gutenberg-tm name associated with
the work.  You can easily comply with the terms of this agreement by
keeping this work in the same format with its attached full Project
Gutenberg-tm License when you share it without charge with others.

1.D.  The copyright laws of the place where you are located also govern
what you can do with this work.  Copyright laws in most countries are in
a constant state of change.  If you are outside the United States, check
the laws of your country in addition to the terms of this agreement
before downloading, copying, displaying, performing, distributing or
creating derivative works based on this work or any other Project
Gutenberg-tm work.  The Foundation makes no representations concerning
the copyright status of any work in any country outside the United
States.

1.E.  Unless you have removed all references to Project Gutenberg:

1.E.1.  The following sentence, with active links to, or other immediate
access to, the full Project Gutenberg-tm License must appear prominently
whenever any copy of a Project Gutenberg-tm work (any work on which the
phrase "Project Gutenberg" appears, or with which the phrase "Project
Gutenberg" is associated) is accessed, displayed, performed, viewed,
copied or distributed:

This eBook is for the use of anyone anywhere at no cost and with
almost no restrictions whatsoever.  You may copy it, give it away or
re-use it under the terms of the Project Gutenberg License included
with this eBook or online at www.gutenberg.org

1.E.2.  If an individual Project Gutenberg-tm electronic work is derived
from the public domain (does not contain a notice indicating that it is
posted with permission of the copyright holder), the work can be copied
and distributed to anyone in the United States without paying any fees
or charges.  If you are redistributing or providing access to a work
with the phrase "Project Gutenberg" associated with or appearing on the
work, you must comply either with the requirements of paragraphs 1.E.1
through 1.E.7 or obtain permission for the use of the work and the
Project Gutenberg-tm trademark as set forth in paragraphs 1.E.8 or
1.E.9.

1.E.3.  If an individual Project Gutenberg-tm electronic work is posted
with the permission of the copyright holder, your use and distribution
must comply with both paragraphs 1.E.1 through 1.E.7 and any additional
terms imposed by the copyright holder.  Additional terms will be linked
to the Project Gutenberg-tm License for all works posted with the
permission of the copyright holder found at the beginning of this work.

1.E.4.  Do not unlink or detach or remove the full Project Gutenberg-tm
License terms from this work, or any files containing a part of this
work or any other work associated with Project Gutenberg-tm.

1.E.5.  Do not copy, display, perform, distribute or redistribute this
electronic work, or any part of this electronic work, without
prominently displaying the sentence set forth in paragraph 1.E.1 with
active links or immediate access to the full terms of the Project
Gutenberg-tm License.

1.E.6.  You may convert to and distribute this work in any binary,
compressed, marked up, nonproprietary or proprietary form, including any
word processing or hypertext form.  However, if you provide access to or
distribute copies of a Project Gutenberg-tm work in a format other than
"Plain Vanilla ASCII" or other format used in the official version
posted on the official Project Gutenberg-tm web site (www.gutenberg.org),
you must, at no additional cost, fee or expense to the user, provide a
copy, a means of exporting a copy, or a means of obtaining a copy upon
request, of the work in its original "Plain Vanilla ASCII" or other
form.  Any alternate format must include the full Project Gutenberg-tm
License as specified in paragraph 1.E.1.

1.E.7.  Do not charge a fee for access to, viewing, displaying,
performing, copying or distributing any Project Gutenberg-tm works
unless you comply with paragraph 1.E.8 or 1.E.9.

1.E.8.  You may charge a reasonable fee for copies of or providing
access to or distributing Project Gutenberg-tm electronic works provided
that

- You pay a royalty fee of 20% of the gross profits you derive from
     the use of Project Gutenberg-tm works calculated using the method
     you already use to calculate your applicable taxes.  The fee is
     owed to the owner of the Project Gutenberg-tm trademark, but he
     has agreed to donate royalties under this paragraph to the
     Project Gutenberg Literary Archive Foundation.  Royalty payments
     must be paid within 60 days following each date on which you
     prepare (or are legally required to prepare) your periodic tax
     returns.  Royalty payments should be clearly marked as such and
     sent to the Project Gutenberg Literary Archive Foundation at the
     address specified in Section 4, "Information about donations to
     the Project Gutenberg Literary Archive Foundation."

- You provide a full refund of any money paid by a user who notifies
     you in writing (or by e-mail) within 30 days of receipt that s/he
     does not agree to the terms of the full Project Gutenberg-tm
     License.  You must require such a user to return or
     destroy all copies of the works possessed in a physical medium
     and discontinue all use of and all access to other copies of
     Project Gutenberg-tm works.

- You provide, in accordance with paragraph 1.F.3, a full refund of any
     money paid for a work or a replacement copy, if a defect in the
     electronic work is discovered and reported to you within 90 days
     of receipt of the work.

- You comply with all other terms of this agreement for free
     distribution of Project Gutenberg-tm works.

1.E.9.  If you wish to charge a fee or distribute a Project Gutenberg-tm
electronic work or group of works on different terms than are set
forth in this agreement, you must obtain permission in writing from
both the Project Gutenberg Literary Archive Foundation and Michael
Hart, the owner of the Project Gutenberg-tm trademark.  Contact the
Foundation as set forth in Section 3 below.

1.F.

1.F.1.  Project Gutenberg volunteers and employees expend considerable
effort to identify, do copyright research on, transcribe and proofread
public domain works in creating the Project Gutenberg-tm
collection.  Despite these efforts, Project Gutenberg-tm electronic
works, and the medium on which they may be stored, may contain
"Defects," such as, but not limited to, incomplete, inaccurate or
corrupt data, transcription errors, a copyright or other intellectual
property infringement, a defective or damaged disk or other medium, a
computer virus, or computer codes that damage or cannot be read by
your equipment.

1.F.2.  LIMITED WARRANTY, DISCLAIMER OF DAMAGES - Except for the "Right
of Replacement or Refund" described in paragraph 1.F.3, the Project
Gutenberg Literary Archive Foundation, the owner of the Project
Gutenberg-tm trademark, and any other party distributing a Project
Gutenberg-tm electronic work under this agreement, disclaim all
liability to you for damages, costs and expenses, including legal
fees.  YOU AGREE THAT YOU HAVE NO REMEDIES FOR NEGLIGENCE, STRICT
LIABILITY, BREACH OF WARRANTY OR BREACH OF CONTRACT EXCEPT THOSE
PROVIDED IN PARAGRAPH F3.  YOU AGREE THAT THE FOUNDATION, THE
TRADEMARK OWNER, AND ANY DISTRIBUTOR UNDER THIS AGREEMENT WILL NOT BE
LIABLE TO YOU FOR ACTUAL, DIRECT, INDIRECT, CONSEQUENTIAL, PUNITIVE OR
INCIDENTAL DAMAGES EVEN IF YOU GIVE NOTICE OF THE POSSIBILITY OF SUCH
DAMAGE.

1.F.3.  LIMITED RIGHT OF REPLACEMENT OR REFUND - If you discover a
defect in this electronic work within 90 days of receiving it, you can
receive a refund of the money (if any) you paid for it by sending a
written explanation to the person you received the work from.  If you
received the work on a physical medium, you must return the medium with
your written explanation.  The person or entity that provided you with
the defective work may elect to provide a replacement copy in lieu of a
refund.  If you received the work electronically, the person or entity
providing it to you may choose to give you a second opportunity to
receive the work electronically in lieu of a refund.  If the second copy
is also defective, you may demand a refund in writing without further
opportunities to fix the problem.

1.F.4.  Except for the limited right of replacement or refund set forth
in paragraph 1.F.3, this work is provided to you 'AS-IS' WITH NO OTHER
WARRANTIES OF ANY KIND, EXPRESS OR IMPLIED, INCLUDING BUT NOT LIMITED TO
WARRANTIES OF MERCHANTIBILITY OR FITNESS FOR ANY PURPOSE.

1.F.5.  Some states do not allow disclaimers of certain implied
warranties or the exclusion or limitation of certain types of damages.
If any disclaimer or limitation set forth in this agreement violates the
law of the state applicable to this agreement, the agreement shall be
interpreted to make the maximum disclaimer or limitation permitted by
the applicable state law.  The invalidity or unenforceability of any
provision of this agreement shall not void the remaining provisions.

1.F.6.  INDEMNITY - You agree to indemnify and hold the Foundation, the
trademark owner, any agent or employee of the Foundation, anyone
providing copies of Project Gutenberg-tm electronic works in accordance
with this agreement, and any volunteers associated with the production,
promotion and distribution of Project Gutenberg-tm electronic works,
harmless from all liability, costs and expenses, including legal fees,
that arise directly or indirectly from any of the following which you do
or cause to occur: (a) distribution of this or any Project Gutenberg-tm
work, (b) alteration, modification, or additions or deletions to any
Project Gutenberg-tm work, and (c) any Defect you cause.


Section  2.  Information about the Mission of Project Gutenberg-tm

Project Gutenberg-tm is synonymous with the free distribution of
electronic works in formats readable by the widest variety of computers
including obsolete, old, middle-aged and new computers.  It exists
because of the efforts of hundreds of volunteers and donations from
people in all walks of life.

Volunteers and financial support to provide volunteers with the
assistance they need, is critical to reaching Project Gutenberg-tm's
goals and ensuring that the Project Gutenberg-tm collection will
remain freely available for generations to come.  In 2001, the Project
Gutenberg Literary Archive Foundation was created to provide a secure
and permanent future for Project Gutenberg-tm and future generations.
To learn more about the Project Gutenberg Literary Archive Foundation
and how your efforts and donations can help, see Sections 3 and 4
and the Foundation web page at http://www.pglaf.org.


Section 3.  Information about the Project Gutenberg Literary Archive
Foundation

The Project Gutenberg Literary Archive Foundation is a non profit
501(c)(3) educational corporation organized under the laws of the
state of Mississippi and granted tax exempt status by the Internal
Revenue Service.  The Foundation's EIN or federal tax identification
number is 64-6221541.  Its 501(c)(3) letter is posted at
http://pglaf.org/fundraising.  Contributions to the Project Gutenberg
Literary Archive Foundation are tax deductible to the full extent
permitted by U.S. federal laws and your state's laws.

The Foundation's principal office is located at 4557 Melan Dr. S.
Fairbanks, AK, 99712., but its volunteers and employees are scattered
throughout numerous locations.  Its business office is located at
809 North 1500 West, Salt Lake City, UT 84116, (801) 596-1887, email
business@pglaf.org.  Email contact links and up to date contact
information can be found at the Foundation's web site and official
page at http://pglaf.org

For additional contact information:
     Dr. Gregory B. Newby
     Chief Executive and Director
     gbnewby@pglaf.org


Section 4.  Information about Donations to the Project Gutenberg
Literary Archive Foundation

Project Gutenberg-tm depends upon and cannot survive without wide
spread public support and donations to carry out its mission of
increasing the number of public domain and licensed works that can be
freely distributed in machine readable form accessible by the widest
array of equipment including outdated equipment.  Many small donations
($1 to $5,000) are particularly important to maintaining tax exempt
status with the IRS.

The Foundation is committed to complying with the laws regulating
charities and charitable donations in all 50 states of the United
States.  Compliance requirements are not uniform and it takes a
considerable effort, much paperwork and many fees to meet and keep up
with these requirements.  We do not solicit donations in locations
where we have not received written confirmation of compliance.  To
SEND DONATIONS or determine the status of compliance for any
particular state visit http://pglaf.org

While we cannot and do not solicit contributions from states where we
have not met the solicitation requirements, we know of no prohibition
against accepting unsolicited donations from donors in such states who
approach us with offers to donate.

International donations are gratefully accepted, but we cannot make
any statements concerning tax treatment of donations received from
outside the United States.  U.S. laws alone swamp our small staff.

Please check the Project Gutenberg Web pages for current donation
methods and addresses.  Donations are accepted in a number of other
ways including checks, online payments and credit card donations.
To donate, please visit: http://pglaf.org/donate


Section 5.  General Information About Project Gutenberg-tm electronic
works.

Professor Michael S. Hart is the originator of the Project Gutenberg-tm
concept of a library of electronic works that could be freely shared
with anyone.  For thirty years, he produced and distributed Project
Gutenberg-tm eBooks with only a loose network of volunteer support.


Project Gutenberg-tm eBooks are often created from several printed
editions, all of which are confirmed as Public Domain in the U.S.
unless a copyright notice is included.  Thus, we do not necessarily
keep eBooks in compliance with any particular paper edition.


Most people start at our Web site which has the main PG search facility:

     http://www.gutenberg.org

This Web site includes information about Project Gutenberg-tm,
including how to make donations to the Project Gutenberg Literary
Archive Foundation, how to help produce our new eBooks, and how to
subscribe to our email newsletter to hear about new eBooks.
\end{PGboilerplate}

% %%%%%%%%%%%%%%%%%%%%%%%%%%%%%%%%%%%%%%%%%%%%%%%%%%%%%%%%%%%%%%%%%%%%%%% %
%                                                                         %
% End of Project Gutenberg's How to Draw a Straight Line, by A.B. Kempe   %
%                                                                         %
% *** END OF THIS PROJECT GUTENBERG EBOOK HOW TO DRAW A STRAIGHT LINE *** %
%                                                                         %
% ***** This file should be named 25155-t.tex or 25155-t.zip *****        %
% This and all associated files of various formats will be found in:      %
%         http://www.gutenberg.org/2/5/1/5/25155/                         %
%                                                                         %
% %%%%%%%%%%%%%%%%%%%%%%%%%%%%%%%%%%%%%%%%%%%%%%%%%%%%%%%%%%%%%%%%%%%%%%% %

\end{document}

### lprep configuration ###
$PageSeparator = qr/^\\PGx?-/;
@ControlwordReplace = (
  ['\\dotsc', '...'],
                      );
@ControlwordArguments = (
  ['\\Note', 1, 1, ' (', ')'],
  ['\\figref', 1, 1, 'Fig.~', ''],
  ['\\ForceBookmark', 1, 0, '', '', 1, 0, '', ''],
                        );

$CustomClean = 'print "\\nCustom cleaning in progress...";
my $cline = 0;
while ($cline <=$#file) {
  $file[$cline] =~ s/\\|.*$//g; # strip comments in PGboilerplate
  $cline++
}
print "done\\n";';

###
This is pdfeTeX, Version 3.141592-1.21a-2.2 (Web2C 7.5.4) (format=pdflatex 2007.10.4)  24 APR 2008 09:07
entering extended mode
**25155-t.tex
(./25155-t.tex
LaTeX2e <2003/12/01>
Babel <v3.8d> and hyphenation patterns for american, french, german, ngerman, b
ahasa, basque, bulgarian, catalan, croatian, czech, danish, dutch, esperanto, e
stonian, finnish, greek, icelandic, irish, italian, latin, magyar, norsk, polis
h, portuges, romanian, russian, serbian, slovak, slovene, spanish, swedish, tur
kish, ukrainian, nohyphenation, loaded.
\PGheader=\toks14
(/usr/share/texmf-tetex/tex/latex/memoir/memoir.cls
Document Class: memoir 2004/04/05 v1.61 configurable document class
\onelineskip=\skip41
\lxvchars=\skip42
\xlvchars=\skip43
\@memcnta=\count79
\stockheight=\skip44
\stockwidth=\skip45
\trimtop=\skip46
\trimedge=\skip47
(/usr/share/texmf-tetex/tex/latex/memoir/mem12.clo
File: mem12.clo 2004/03/12 v0.3 memoir class 12pt size option
)
\spinemargin=\skip48
\foremargin=\skip49
\uppermargin=\skip50
\lowermargin=\skip51
\headdrop=\skip52
\normalrulethickness=\skip53
\headwidth=\skip54
\c@storedpagenumber=\count80
\thanksmarkwidth=\skip55
\thanksmarksep=\skip56
\droptitle=\skip57
\abstitleskip=\skip58
\absleftindent=\skip59
\absrightindent=\skip60
\absparindent=\skip61
\absparsep=\skip62
\c@part=\count81
\c@chapter=\count82
\c@section=\count83
\c@subsection=\count84
\c@subsubsection=\count85
\c@paragraph=\count86
\c@subparagraph=\count87
\beforechapskip=\skip63
\midchapskip=\skip64
\afterchapskip=\skip65
\chapindent=\skip66
\bottomsectionskip=\skip67
\secindent=\skip68
\beforesecskip=\skip69
\aftersecskip=\skip70
\subsecindent=\skip71
\beforesubsecskip=\skip72
\aftersubsecskip=\skip73
\subsubsecindent=\skip74
\beforesubsubsecskip=\skip75
\aftersubsubsecskip=\skip76
\paraindent=\skip77
\beforeparaskip=\skip78
\afterparaskip=\skip79
\subparaindent=\skip80
\beforesubparaskip=\skip81
\aftersubparaskip=\skip82
\pfbreakskip=\skip83
\c@@ppsavesec=\count88
\c@@ppsaveapp=\count89
\ragrparindent=\dimen102
\parsepi=\skip84
\topsepi=\skip85
\itemsepi=\skip86
\parsepii=\skip87
\topsepii=\skip88
\topsepiii=\skip89
\m@msavetopsep=\skip90
\m@msavepartopsep=\skip91
\@enLab=\toks15
\c@vslineno=\count90
\c@poemline=\count91
\c@modulo@vs=\count92
\vleftskip=\skip92
\vrightskip=\skip93
\stanzaskip=\skip94
\versewidth=\skip95
\vgap=\skip96
\vindent=\skip97
\c@verse=\count93
\c@chrsinstr=\count94
\beforepoemtitleskip=\skip98
\afterpoemtitleskip=\skip99
\col@sep=\dimen103
\extrarowheight=\dimen104
\NC@list=\toks16
\extratabsurround=\skip100
\backup@length=\skip101
\TX@col@width=\dimen105
\TX@old@table=\dimen106
\TX@old@col=\dimen107
\TX@target=\dimen108
\TX@delta=\dimen109
\TX@cols=\count95
\TX@ftn=\toks17
\heavyrulewidth=\dimen110
\lightrulewidth=\dimen111
\cmidrulewidth=\dimen112
\belowrulesep=\dimen113
\belowbottomsep=\dimen114
\aboverulesep=\dimen115
\abovetopsep=\dimen116
\cmidrulesep=\dimen117
\cmidrulekern=\dimen118
\defaultaddspace=\dimen119
\@cmidla=\count96
\@cmidlb=\count97
\@aboverulesep=\dimen120
\@belowrulesep=\dimen121
\@thisruleclass=\count98
\@lastruleclass=\count99
\@thisrulewidth=\dimen122
\ctableftskip=\skip102
\ctabrightskip=\skip103
\abovecolumnspenalty=\count100
\@linestogo=\count101
\@cellstogo=\count102
\@cellsincolumn=\count103
\crtok=\toks18
\@mincolumnwidth=\dimen123
\c@newflo@tctr=\count104
\@contcwidth=\skip104
\@contindw=\skip105
\abovecaptionskip=\skip106
\belowcaptionskip=\skip107
\subfloattopskip=\skip108
\subfloatcapskip=\skip109
\subfloatcaptopadj=\skip110
\subfloatbottomskip=\skip111
\subfloatlabelskip=\skip112
\subfloatcapmargin=\dimen124
\c@@contsubnum=\count105
\beforeepigraphskip=\skip113
\afterepigraphskip=\skip114
\epigraphwidth=\skip115
\epigraphrule=\skip116
LaTeX Info: Redefining \emph on input line 4419.
LaTeX Info: Redefining \em on input line 4420.
\tocentryskip=\skip117
\tocbaseline=\skip118
\cftparskip=\skip119
\cftbeforepartskip=\skip120
\cftpartindent=\skip121
\cftpartnumwidth=\skip122
\cftbeforechapterskip=\skip123
\cftchapterindent=\skip124
\cftchapternumwidth=\skip125
\cftbeforesectionskip=\skip126
\cftsectionindent=\skip127
\cftsectionnumwidth=\skip128
\cftbeforesubsectionskip=\skip129
\cftsubsectionindent=\skip130
\cftsubsectionnumwidth=\skip131
\cftbeforesubsubsectionskip=\skip132
\cftsubsubsectionindent=\skip133
\cftsubsubsectionnumwidth=\skip134
\cftbeforeparagraphskip=\skip135
\cftparagraphindent=\skip136
\cftparagraphnumwidth=\skip137
\cftbeforesubparagraphskip=\skip138
\cftsubparagraphindent=\skip139
\cftsubparagraphnumwidth=\skip140
\c@maxsecnumdepth=\count106
\bibindent=\dimen125
\bibitemsep=\skip141
\indexcolsep=\skip142
\indexrule=\skip143
\indexmarkstyle=\toks19
\@indexbox=\insert233
\sideparvshift=\skip144
\sideins=\insert232
\sidebarhsep=\skip145
\sidebarvsep=\skip146
\sidebarwidth=\skip147
\footmarkwidth=\skip148
\footmarksep=\skip149
\footparindent=\skip150
\footinsdim=\skip151
\footinsv@r=\insert231
\@mpfootinsv@r=\insert230
\m@m@k=\count107
\m@m@h=\dimen126
\m@mipn@skip=\skip152
\c@sheetsequence=\count108
\c@lastsheet=\count109
\c@lastpage=\count110
\every@verbatim=\toks20
\afterevery@verbatim=\toks21
\verbatim@line=\toks22
\tab@position=\count111
\verbatim@in@stream=\read1
\verbatimindent=\skip153
\verbatim@out=\write3
\bvboxsep=\skip154
\c@bvlinectr=\count112
\bvnumlength=\skip155
\FrameRule=\dimen127
\FrameSep=\dimen128
\c@cp@cntr=\count113
LaTeX Info: Redefining \: on input line 7541.
LaTeX Info: Redefining \! on input line 7543.
\c@ism@mctr=\count114
\c@xsm@mctr=\count115
\c@csm@mctr=\count116
\c@ksm@mctr=\count117
\c@xksm@mctr=\count118
\c@cksm@mctr=\count119
\c@msm@mctr=\count120
\c@xmsm@mctr=\count121
\c@cmsm@mctr=\count122
\c@bsm@mctr=\count123
\c@workm@mctr=\count124
\c@figure=\count125
\c@lofdepth=\count126
\c@lofdepth=\count126
\cftbeforefigureskip=\skip156
\cftfigureindent=\skip157
\cftfigurenumwidth=\skip158
\c@table=\count126
\c@lotdepth=\count127
\c@lotdepth=\count127
\cftbeforetableskip=\skip159
\cfttableindent=\skip160
\cfttablenumwidth=\skip161
(/usr/share/texmf-tetex/tex/latex/memoir/mempatch.sty
File: mempatch.sty 2005/02/01 v3.5 Patches for memoir class v1.61
\abs@leftindent=\dimen129
LaTeX Info: Redefining \em on input line 142.
LaTeX Info: Redefining \em on input line 384.
LaTeX Info: Redefining \emph on input line 392.
))

LaTeX Warning: You have requested, on input line 111, version
               `2005/09/25' of document class memoir,
               but only version
               `2004/04/05 v1.61 configurable document class'
               is available.


******************************************************
Stock height and width: 654.41338pt by 423.94606pt
Top and edge trims: 0.0pt and 0.0pt
Page height and width: 654.41338pt by 423.94606pt
Text height and width: 523.53055pt by 312.0pt
Spine and edge margins: 51.21504pt and 59.75095pt
Upper and lower margins: 59.75095pt and 71.13188pt
Headheight and headsep: 14.0pt and 14.0pt
Footskip: 17.07182pt
Columnsep and columnseprule: 10.0pt and 0.0pt
Marginparsep and marginparwidth: 7.0pt and 102.0pt
******************************************************

LaTeX Info: Redefining \ttfamily on input line 130.
(/usr/share/texmf-tetex/tex/latex/base/inputenc.sty
Package: inputenc 2004/02/05 v1.0d Input encoding file
(/usr/share/texmf-tetex/tex/latex/base/latin1.def
File: latin1.def 2004/02/05 v1.0d Input encoding file
))

LaTeX Warning: You have requested, on input line 140, version
               `2006/05/05' of package inputenc,
               but only version
               `2004/02/05 v1.0d Input encoding file'
               is available.

(/usr/share/texmf-tetex/tex/latex/amsmath/amsmath.sty
Package: amsmath 2000/07/18 v2.13 AMS math features
\@mathmargin=\skip162
For additional information on amsmath, use the `?' option.
(/usr/share/texmf-tetex/tex/latex/amsmath/amstext.sty
Package: amstext 2000/06/29 v2.01
(/usr/share/texmf-tetex/tex/latex/amsmath/amsgen.sty
File: amsgen.sty 1999/11/30 v2.0
\@emptytoks=\toks23
\ex@=\dimen130
)) (/usr/share/texmf-tetex/tex/latex/amsmath/amsbsy.sty
Package: amsbsy 1999/11/29 v1.2d
\pmbraise@=\dimen131
) (/usr/share/texmf-tetex/tex/latex/amsmath/amsopn.sty
Package: amsopn 1999/12/14 v2.01 operator names
)
\inf@bad=\count127
LaTeX Info: Redefining \frac on input line 211.
\uproot@=\count128
\leftroot@=\count129
LaTeX Info: Redefining \overline on input line 307.
\classnum@=\count130
\DOTSCASE@=\count131
LaTeX Info: Redefining \ldots on input line 379.
LaTeX Info: Redefining \dots on input line 382.
LaTeX Info: Redefining \cdots on input line 467.
\Mathstrutbox@=\box26
\strutbox@=\box27
\big@size=\dimen132
LaTeX Font Info:    Redeclaring font encoding OML on input line 567.
LaTeX Font Info:    Redeclaring font encoding OMS on input line 568.
\macc@depth=\count132
\c@MaxMatrixCols=\count133
\dotsspace@=\muskip10
\c@parentequation=\count134
\dspbrk@lvl=\count135
\tag@help=\toks24
\row@=\count136
\column@=\count137
\maxfields@=\count138
\andhelp@=\toks25
\eqnshift@=\dimen133
\alignsep@=\dimen134
\tagshift@=\dimen135
\tagwidth@=\dimen136
\totwidth@=\dimen137
\lineht@=\dimen138
\@envbody=\toks26
\multlinegap=\skip163
\multlinetaggap=\skip164
\mathdisplay@stack=\toks27
LaTeX Info: Redefining \[ on input line 2666.
LaTeX Info: Redefining \] on input line 2667.
) (/usr/share/texmf-tetex/tex/latex/soul/soul.sty
Package: soul 2003/11/17 v2.4 letterspacing/underlining (mf)
\SOUL@word=\toks28
\SOUL@lasttoken=\toks29
\SOUL@cmds=\toks30
\SOUL@buffer=\toks31
\SOUL@token=\toks32
\SOUL@spaceskip=\skip165
\SOUL@ttwidth=\dimen139
\SOUL@uldp=\dimen140
\SOUL@ulht=\dimen141
) (/usr/share/texmf-tetex/tex/latex/tools/indentfirst.sty
Package: indentfirst 1995/11/23 v1.03 Indent first paragraph (DPC)
) (/usr/share/texmf-tetex/tex/latex/graphics/graphicx.sty
Package: graphicx 1999/02/16 v1.0f Enhanced LaTeX Graphics (DPC,SPQR)
(/usr/share/texmf-tetex/tex/latex/graphics/keyval.sty
Package: keyval 1999/03/16 v1.13 key=value parser (DPC)
\KV@toks@=\toks33
) (/usr/share/texmf-tetex/tex/latex/graphics/graphics.sty
Package: graphics 2001/07/07 v1.0n Standard LaTeX Graphics (DPC,SPQR)
(/usr/share/texmf-tetex/tex/latex/graphics/trig.sty
Package: trig 1999/03/16 v1.09 sin cos tan (DPC)
) (/usr/share/texmf-tetex/tex/latex/graphics/graphics.cfg
File: graphics.cfg 2005/02/03 v1.3 graphics configuration of teTeX/TeXLive
)
Package graphics Info: Driver file: pdftex.def on input line 80.
(/usr/share/texmf-tetex/tex/latex/graphics/pdftex.def
File: pdftex.def 2002/06/19 v0.03k graphics/color for pdftex
\Gread@gobject=\count139
))
\Gin@req@height=\dimen142
\Gin@req@width=\dimen143
)

***
*** Important Note: this document comes with PNG and PDF graphics,
*** so make sure you use an appropriate workflow!
*** 

(/usr/share/texmf-tetex/tex/latex/wrapfig/wrapfig.sty
\wrapoverhang=\dimen144
\WF@size=\dimen145
\c@WF@wrappedlines=\count140
\WF@box=\box28
\WF@everypar=\toks34
Package: wrapfig 2003/01/31  v 3.6
) (/usr/share/texmf-tetex/tex/latex/hyperref/hyperref.sty
Package: hyperref 2003/11/30 v6.74m Hypertext links for LaTeX
\@linkdim=\dimen146
\Hy@linkcounter=\count141
\Hy@pagecounter=\count142
(/usr/share/texmf-tetex/tex/latex/hyperref/pd1enc.def
File: pd1enc.def 2003/11/30 v6.74m Hyperref: PDFDocEncoding definition (HO)
) (/usr/share/texmf-tetex/tex/latex/hyperref/hyperref.cfg
File: hyperref.cfg 2002/06/06 v1.2 hyperref configuration of TeXLive and teTeX
)
Package hyperref Info: Option `draft' set `true' on input line 1830.
Package hyperref Info: Option `colorlinks' set `true' on input line 1830.
Package hyperref Info: Hyper figures OFF on input line 1880.
Package hyperref Info: Link nesting OFF on input line 1885.
Package hyperref Info: Hyper index ON on input line 1888.
Package hyperref Info: Plain pages ON on input line 1893.
Package hyperref Info: Backreferencing OFF on input line 1900.
Implicit mode ON; LaTeX internals redefined
Package hyperref Info: Bookmarks ON on input line 2004.
(/usr/share/texmf-tetex/tex/latex/url/url.sty
\Urlmuskip=\muskip11
Package: url 2004/03/15  ver 3.1  Verb mode for urls, etc.
)
LaTeX Info: Redefining \url on input line 2143.
\Fld@menulength=\count143
\Field@Width=\dimen147
\Fld@charsize=\dimen148
\Choice@toks=\toks35
\Field@toks=\toks36
Package hyperref Info: Hyper figures OFF on input line 2618.
Package hyperref Info: Link nesting OFF on input line 2623.
Package hyperref Info: Hyper index ON on input line 2626.
Package hyperref Info: backreferencing OFF on input line 2633.
Package hyperref Info: Link coloring ON on input line 2636.
\c@Item=\count144
)
*hyperref using default driver hpdftex*
(/usr/share/texmf-tetex/tex/latex/hyperref/hpdftex.def
File: hpdftex.def 2003/11/30 v6.74m Hyperref driver for pdfTeX
(/usr/share/texmf-tetex/tex/latex/psnfss/pifont.sty
Package: pifont 2004/09/15 PSNFSS-v9.2 Pi font support (SPQR) 
LaTeX Font Info:    Try loading font information for U+pzd on input line 63.
(/usr/share/texmf-tetex/tex/latex/psnfss/upzd.fd
File: upzd.fd 2001/06/04 font definitions for U/pzd.
)
LaTeX Font Info:    Try loading font information for U+psy on input line 64.
(/usr/share/texmf-tetex/tex/latex/psnfss/upsy.fd
File: upsy.fd 2001/06/04 font definitions for U/psy.
))
\Fld@listcount=\count145
\@outlinefile=\write4
) (/usr/share/texmf-tetex/tex/latex/memoir/memhfixc.sty
Package: memhfixc 2004/05/13 v1.6 package fixes for memoir class
)
Package hyperref Info: Option `pdfdisplaydoctitle' set `true' on input line 300
.
Package hyperref Info: Option `bookmarksopen' set `true' on input line 300.
Package hyperref Info: Option `linktocpage' set `false' on input line 300.
Package hyperref Info: Option `plainpages' set `false' on input line 300.
(./25155-t.aux)
\openout1 = `25155-t.aux'.

LaTeX Font Info:    Checking defaults for OML/cmm/m/it on input line 373.
LaTeX Font Info:    ... okay on input line 373.
LaTeX Font Info:    Checking defaults for T1/cmr/m/n on input line 373.
LaTeX Font Info:    ... okay on input line 373.
LaTeX Font Info:    Checking defaults for OT1/cmr/m/n on input line 373.
LaTeX Font Info:    ... okay on input line 373.
LaTeX Font Info:    Checking defaults for OMS/cmsy/m/n on input line 373.
LaTeX Font Info:    ... okay on input line 373.
LaTeX Font Info:    Checking defaults for OMX/cmex/m/n on input line 373.
LaTeX Font Info:    ... okay on input line 373.
LaTeX Font Info:    Checking defaults for U/cmr/m/n on input line 373.
LaTeX Font Info:    ... okay on input line 373.
LaTeX Font Info:    Checking defaults for PD1/pdf/m/n on input line 373.
LaTeX Font Info:    ... okay on input line 373.
\c@lofdepth=\count146
\c@lotdepth=\count147
(/usr/share/texmf-tetex/tex/context/base/supp-pdf.tex (/usr/share/texmf-tetex/t
ex/context/base/supp-mis.tex
loading : Context Support Macros / Miscellaneous (2004.10.26)
\protectiondepth=\count148
\scratchcounter=\count149
\scratchtoks=\toks37
\scratchdimen=\dimen149
\scratchskip=\skip166
\scratchmuskip=\muskip12
\scratchbox=\box29
\scratchread=\read2
\scratchwrite=\write5
\zeropoint=\dimen150
\onepoint=\dimen151
\onebasepoint=\dimen152
\minusone=\count150
\thousandpoint=\dimen153
\onerealpoint=\dimen154
\emptytoks=\toks38
\nextbox=\box30
\nextdepth=\dimen155
\everyline=\toks39
\!!counta=\count151
\!!countb=\count152
\recursecounter=\count153
)
loading : Context Support Macros / PDF (2004.03.26)
\nofMPsegments=\count154
\nofMParguments=\count155
\MPscratchCnt=\count156
\MPscratchDim=\dimen156
\MPnumerator=\count157
\everyMPtoPDFconversion=\toks40
) (/usr/share/texmf-tetex/tex/latex/graphics/color.sty
Package: color 1999/02/16 v1.0i Standard LaTeX Color (DPC)
LaTeX Info: Redefining \color on input line 71.
(/usr/share/texmf-tetex/tex/latex/graphics/color.cfg
File: color.cfg 2005/02/03 v1.3 color configuration of teTeX/TeXLive
)
Package color Info: Driver file: pdftex.def on input line 125.
)
Package hyperref Info: Link coloring ON on input line 373.
(/usr/share/texmf-tetex/tex/latex/hyperref/nameref.sty
Package: nameref 2003/12/03 v2.21 Cross-referencing by name of section
\c@section@level=\count158
)
LaTeX Info: Redefining \ref on input line 373.
LaTeX Info: Redefining \pageref on input line 373.
(./25155-t.out) (./25155-t.out)
\openout4 = `25155-t.out'.

Redoing nameref's sectioning
Redoing nameref's label
LaTeX Font Info:    Try loading font information for T1+pcr on input line 377.
(/usr/share/texmf-tetex/tex/latex/psnfss/t1pcr.fd
File: t1pcr.fd 2001/06/04 font definitions for T1/pcr.
) [1


{/var/lib/texmf/fonts/map/pdftex/updmap/pdftex.map}] [2

] <images/illus0003.png, id=176, 105.5142pt x 44.5665pt>
File: images/illus0003.png Graphic file (type png)
<use images/illus0003.png> [3 <./images/illus0003.png>] [4

] <images/london.pdf, id=191, 182.6825pt x 72.27pt>
File: images/london.pdf Graphic file (type pdf)
<use images/london.pdf> [5 <./images/london.pdf>] [6

] [7


] [8

]
Underfull \vbox (badness 2644) has occurred while \output is active []

[1

] [2]
Underfull \vbox (badness 10000) has occurred while \output is active []

[3]
Underfull \vbox (badness 10000) has occurred while \output is active []

[4] <images/illus0015.png, id=276, 115.1502pt x 56.6115pt>
File: images/illus0015.png Graphic file (type png)
<use images/illus0015.png> [5 <./images/illus0015.png>]
Underfull \vbox (badness 10000) has occurred while \output is active []

[6]
Underfull \vbox (badness 10000) has occurred while \output is active []

[7] <images/illus0019.png, id=293, 140.9265pt x 81.1833pt>
File: images/illus0019.png Graphic file (type png)
<use images/illus0019.png> [8 <./images/illus0019.png>] <images/illus0020.png, 
id=302, 77.5698pt x 98.0463pt>
File: images/illus0020.png Graphic file (type png)
<use images/illus0020.png>
Underfull \vbox (badness 5316) has occurred while \output is active []

[9] <images/illus0021.png, id=308, 78.0516pt x 97.8054pt>
File: images/illus0021.png Graphic file (type png)
<use images/illus0021.png> [10 <./images/illus0020.png> <./images/illus0021.png
>] <images/illus0022.png, id=314, 116.1138pt x 70.3428pt>
File: images/illus0022.png Graphic file (type png)
<use images/illus0022.png> [11 <./images/illus0022.png>] <images/illus0024.png,
 id=320, 177.7842pt x 129.1224pt>
File: images/illus0024.png Graphic file (type png)
<use images/illus0024.png>
Underfull \vbox (badness 6364) has occurred while \output is active []

[12] <images/illus0025.png, id=333, 132.7359pt x 61.6704pt>
File: images/illus0025.png Graphic file (type png)
<use images/illus0025.png> [13 <./images/illus0024.png>] [14 <./images/illus002
5.png>] <images/illus0026.png, id=344, 152.4897pt x 90.3375pt>
File: images/illus0026.png Graphic file (type png)
<use images/illus0026.png> <images/illus0027a.png, id=345, 125.5089pt x 68.4156
pt>
File: images/illus0027a.png Graphic file (type png)
<use images/illus0027a.png> [15 <./images/illus0026.png>] <images/illus0027b.pn
g, id=355, 100.2144pt x 64.0794pt>
File: images/illus0027b.png Graphic file (type png)
<use images/illus0027b.png> [16 <./images/illus0027a.png> <./images/illus0027b.
png>] <images/illus0029a.png, id=362, 134.904pt x 65.5248pt>
File: images/illus0029a.png Graphic file (type png)
<use images/illus0029a.png> <images/illus0029b.png, id=363, 172.0026pt x 63.838
5pt>
File: images/illus0029b.png Graphic file (type png)
<use images/illus0029b.png> [17 <./images/illus0029a.png>] [18 <./images/illus0
029b.png>] <images/illus0031.png, id=381, 129.8451pt x 119.0046pt>
File: images/illus0031.png Graphic file (type png)
<use images/illus0031.png> [19 <./images/illus0031.png>] <images/illus0032.png,
 id=387, 114.4275pt x 105.0324pt>
File: images/illus0032.png Graphic file (type png)
<use images/illus0032.png> [20] <images/illus0034.png, id=397, 165.7392pt x 195
.6108pt>
File: images/illus0034.png Graphic file (type png)
<use images/illus0034.png> [21 <./images/illus0032.png>] [22 <./images/illus003
4.png>] <images/illus0035.png, id=410, 116.8365pt x 101.6598pt>
File: images/illus0035.png Graphic file (type png)
<use images/illus0035.png> [23 <./images/illus0035.png>]
Underfull \vbox (badness 10000) has occurred while \output is active []

[24] <images/illus0038.png, id=426, 90.0966pt x 217.2918pt>
File: images/illus0038.png Graphic file (type png)
<use images/illus0038.png>
Underfull \vbox (badness 10000) has occurred while \output is active []

[25] [26 <./images/illus0038.png>] <images/illus0039.png, id=439, 84.5559pt x 5
4.6843pt>
File: images/illus0039.png Graphic file (type png)
<use images/illus0039.png> <images/illus0040.png, id=440, 95.6373pt x 61.1886pt
>
File: images/illus0040.png Graphic file (type png)
<use images/illus0040.png> [27 <./images/illus0039.png>] <images/illus0041.png,
 id=450, 120.45pt x 134.4222pt>
File: images/illus0041.png Graphic file (type png)
<use images/illus0041.png> [28 <./images/illus0040.png> <./images/illus0041.png
>] <images/illus0042a.png, id=456, 172.2435pt x 51.0708pt>
File: images/illus0042a.png Graphic file (type png)
<use images/illus0042a.png> <images/illus0042b.png, id=457, 101.9007pt x 146.94
9pt>
File: images/illus0042b.png Graphic file (type png)
<use images/illus0042b.png> <images/illus0044.png, id=458, 108.8868pt x 291.489
pt>
File: images/illus0044.png Graphic file (type png)
<use images/illus0044.png> [29 <./images/illus0042a.png>] <images/illus0043.png
, id=468, 56.1297pt x 84.7968pt>
File: images/illus0043.png Graphic file (type png)
<use images/illus0043.png>
Underfull \hbox (badness 1521) in paragraph at lines 1410--1425
[][]\OT1/cmr/m/n/12 In the linkwork of [][]Fig. 22[][],
 []

[30 <./images/illus0042b.png> <./images/illus0043.png>] [31 <./images/illus0044
.png>] <images/illus0045.png, id=480, 147.9126pt x 166.221pt>
File: images/illus0045.png Graphic file (type png)
<use images/illus0045.png> [32 <./images/illus0045.png>] <images/illus0046.png,
 id=489, 133.6995pt x 102.6234pt>
File: images/illus0046.png Graphic file (type png)
<use images/illus0046.png> <images/illus0047.png, id=490, 156.3441pt x 120.2091
pt>
File: images/illus0047.png Graphic file (type png)
<use images/illus0047.png> [33 <./images/illus0046.png>] <images/illus0048.png,
 id=499, 107.4414pt x 187.1793pt>
File: images/illus0048.png Graphic file (type png)
<use images/illus0048.png> <images/illus0049.png, id=500, 63.5976pt x 197.538pt
>
File: images/illus0049.png Graphic file (type png)
<use images/illus0049.png> [34 <./images/illus0047.png>] [35 <./images/illus004
8.png>] <images/illus0050.png, id=510, 127.4361pt x 98.2872pt>
File: images/illus0050.png Graphic file (type png)
<use images/illus0050.png> [36 <./images/illus0049.png>] <images/illus0051.png,
 id=516, 88.8921pt x 197.7789pt>
File: images/illus0051.png Graphic file (type png)
<use images/illus0051.png> <images/illus0052.png, id=517, 144.0582pt x 219.7008
pt>
File: images/illus0052.png Graphic file (type png)
<use images/illus0052.png> <images/illus0053.png, id=518, 144.0582pt x 227.8914
pt>
File: images/illus0053.png Graphic file (type png)
<use images/illus0053.png> [37 <./images/illus0050.png>] [38 <./images/illus005
1.png>] [39 <./images/illus0052.png>] [40 <./images/illus0053.png>] [41] [42

] [43] [44]
Underfull \vbox (badness 3792) has occurred while \output is active []

[45] [46] <images/illus0060.png, id=598, 62.634pt x 68.6565pt>
File: images/illus0060.png Graphic file (type png)
<use images/illus0060.png> [47 <./images/illus0060.png>] [48

] [49] [50] [51] [52] [53] [54] [55] (./25155-t.aux)

 *File List*
  memoir.cls    2004/04/05 v1.61 configurable document class
   mem12.clo    2004/03/12 v0.3 memoir class 12pt size option
mempatch.sty    2005/02/01 v3.5 Patches for memoir class v1.61
inputenc.sty    2004/02/05 v1.0d Input encoding file
  latin1.def    2004/02/05 v1.0d Input encoding file
 amsmath.sty    2000/07/18 v2.13 AMS math features
 amstext.sty    2000/06/29 v2.01
  amsgen.sty    1999/11/30 v2.0
  amsbsy.sty    1999/11/29 v1.2d
  amsopn.sty    1999/12/14 v2.01 operator names
    soul.sty    2003/11/17 v2.4 letterspacing/underlining (mf)
indentfirst.sty    1995/11/23 v1.03 Indent first paragraph (DPC)
graphicx.sty    1999/02/16 v1.0f Enhanced LaTeX Graphics (DPC,SPQR)
  keyval.sty    1999/03/16 v1.13 key=value parser (DPC)
graphics.sty    2001/07/07 v1.0n Standard LaTeX Graphics (DPC,SPQR)
    trig.sty    1999/03/16 v1.09 sin cos tan (DPC)
graphics.cfg    2005/02/03 v1.3 graphics configuration of teTeX/TeXLive
  pdftex.def    2002/06/19 v0.03k graphics/color for pdftex
 wrapfig.sty    2003/01/31  v 3.6
hyperref.sty    2003/11/30 v6.74m Hypertext links for LaTeX
  pd1enc.def    2003/11/30 v6.74m Hyperref: PDFDocEncoding definition (HO)
hyperref.cfg    2002/06/06 v1.2 hyperref configuration of TeXLive and teTeX
     url.sty    2004/03/15  ver 3.1  Verb mode for urls, etc.
 hpdftex.def    2003/11/30 v6.74m Hyperref driver for pdfTeX
  pifont.sty    2004/09/15 PSNFSS-v9.2 Pi font support (SPQR) 
    upzd.fd    2001/06/04 font definitions for U/pzd.
    upsy.fd    2001/06/04 font definitions for U/psy.
memhfixc.sty    2004/05/13 v1.6 package fixes for memoir class
supp-pdf.tex
   color.sty    1999/02/16 v1.0i Standard LaTeX Color (DPC)
   color.cfg    2005/02/03 v1.3 color configuration of teTeX/TeXLive
 nameref.sty    2003/12/03 v2.21 Cross-referencing by name of section
 25155-t.out
 25155-t.out
   t1pcr.fd    2001/06/04 font definitions for T1/pcr.
images/illus0003.png
images/london.pdf
images/illus0015.png
images/illus0019.png
images/illus0020.png
images/illus0021.png
images/illus0022.png
images/illus0024.png
images/illus0025.png
images/illus0026.png
images/illus0027a.png
images/illus0027b.png
images/illus0029a.png
images/illus0029b.png
images/illus0031.png
images/illus0032.png
images/illus0034.png
images/illus0035.png
images/illus0038.png
images/illus0039.png
images/illus0040.png
images/illus0041.png
images/illus0042a.png
images/illus0042b.png
images/illus0044.png
images/illus0043.png
images/illus0045.png
images/illus0046.png
images/illus0047.png
images/illus0048.png
images/illus0049.png
images/illus0050.png
images/illus0051.png
images/illus0052.png
images/illus0053.png
images/illus0060.png
 ***********

 ) 
Here is how much of TeX's memory you used:
 6001 strings out of 94500
 78479 string characters out of 1175771
 168568 words of memory out of 1000000
 8948 multiletter control sequences out of 10000+50000
 15564 words of font info for 54 fonts, out of 500000 for 2000
 580 hyphenation exceptions out of 8191
 26i,12n,30p,235b,382s stack positions out of 1500i,500n,5000p,200000b,5000s
PDF statistics:
 637 PDF objects out of 300000
 139 named destinations out of 131072
 509 words of extra memory for PDF output out of 65536
</usr/share/texmf-tetex/fonts/type1/bluesky/cm/cmmi10.pfb></usr/share/texmf-t
etex/fonts/type1/bluesky/cm/cmsy8.pfb></usr/share/texmf-tetex/fonts/type1/blues
ky/cm/cmmi12.pfb></usr/share/texmf-tetex/fonts/type1/bluesky/cm/cmsl10.pfb></us
r/share/texmf-tetex/fonts/type1/bluesky/cm/cmcsc10.pfb></usr/share/texmf-tetex/
fonts/type1/bluesky/cm/cmr7.pfb></usr/share/texmf-tetex/fonts/type1/bluesky/cm/
cmti10.pfb></usr/share/texmf-tetex/fonts/type1/bluesky/cm/cmr10.pfb></usr/share
/texmf-tetex/fonts/type1/bluesky/cm/cmr9.pfb></usr/share/texmf-tetex/fonts/type
1/bluesky/cm/cmr6.pfb></usr/share/texmf-tetex/fonts/type1/bluesky/cm/cmdunh10.p
fb></usr/share/texmf-tetex/fonts/type1/bluesky/cm/cmsy10.pfb></usr/share/texmf-
tetex/fonts/type1/bluesky/cm/cmti12.pfb></usr/share/texmf-tetex/fonts/type1/blu
esky/cm/cmr8.pfb></usr/share/texmf-tetex/fonts/type1/bluesky/cm/cmr12.pfb>{/usr
/share/texmf-tetex/fonts/enc/dvips/psnfss/8r.enc}</usr/share/texmf-tetex/fonts/
type1/urw/courier/ucrr8a.pfb>
Output written on 25155-t.pdf (63 pages, 1322151 bytes).
